% =====================================================================
%  Harald Høffding, REMARKS ON PLATO'S DIALOGUE PARMENIDES.  The Royal Danish Academy of Sciences and Letters,
%  Philosophical Communications I, 2.  Copenhagen: Andr. Fred. Høst & Søn; Bianco Lunos Bogtrykkeri, 1920.  65 pp.
%  (Bemærkninger om den platoniske Dialog Parmenides.  Det Kgl. Danske Videnskabernes Selskab, Filosofiske Meddelelser I, 2.)
%  ENGLISH TRANSLATION, from transcription.tex (Danish).
%  GENERATED by ebook-method/platons-parmenides/tr_build.py --write from this template and the four translators' parts in
%  ebook-method/platons-parmenides/tr/ (TRANSLATOR-BRIEF.md; terms: GLOSSARY.md, the title and three key terms approved by
%  Hans 2026-10-20).  Conventions: TRANSLATION-PLAYBOOK.md.  Page numbers of the 1920 print in the margin.
% =====================================================================
\documentclass[12pt,a4paper,openany]{book}
\usepackage[utf8]{inputenc}
\usepackage[T1]{fontenc}
\usepackage{libertinus}
\usepackage{libertinust1math}
\usepackage{textalpha}
\usepackage[english]{babel}
\usepackage[protrusion=true,expansion=true]{microtype}
\usepackage{setspace}
\usepackage[margin=1.3in]{geometry}
\usepackage{fancyhdr}
\usepackage{hyperref}
\hypersetup{pdftitle={Remarks on Plato's Dialogue Parmenides}, pdfauthor={Harald Høffding}, hidelinks}
\setlength{\headheight}{15pt}
\pagestyle{fancy}
\fancyhf{}
\fancyhead[LE,RO]{\thepage}
\fancyhead[C]{\textit{Høffding: Remarks on Plato's Dialogue Parmenides (1920)}}
\setstretch{1.3}
\usepackage{marginnote}
\newcommand{\tocline}[3]{\noindent\makebox[2.2em][r]{#1}\hspace{0.6em}#2\dotfill\ #3\par}
\newcommand{\opage}[1]{\marginnote{\footnotesize\textit{[#1]}}}
\newcommand{\chaphead}[1]{\par\bigskip\phantomsection\begin{center}\bfseries #1\end{center}\medskip\ignorespaces}

\begin{document}

% PDF 3, the title page
\begin{titlepage}
\begin{center}
The Royal Danish Academy of Sciences and Letters.\\[0.2em]
{\small Philosophical Communications \textbf{I}, 2.}\\[0.3em]
\rule{\textwidth}{0.4pt}\\[5em]
{\LARGE REMARKS}\\[0.4em]
{\LARGE ON PLATO'S DIALOGUE}\\[0.4em]
{\LARGE PARMENIDES}\\[2em]
{\small BY}\\[2em]
{\large HARALD HØFFDING}\\[4em]
[the Society's seal]\\[5em]
{\Large COPENHAGEN}\\[0.6em]
{\small\textsc{Chief Agent:}} ANDR. FRED. HØST \& SØN, {\small\textsc{Royal Court Booksellers}}\\[1.2em]
BIANCO LUNO PRINTING HOUSE\\[0.3em]
1920\\[3em]
{\small Translated from the Danish}
\end{center}
\end{titlepage}
\clearpage

% ===== TEXT (pp. 3--65) =====
% --- p. 3 ---
\setcounter{footnote}{0}%
\opage{3}P\textsc{lato}’s \textit{Parmenides} is one of the most disputed of
his dialogues. It falls into two sections, each of which
presents difficulties of its own for Platonic scholarship.

In the first section Parmenides of Elea develops, before
Socrates, who is here portrayed as a young man, various
objections that seem to strike the philosophy,
the so-called doctrine of Ideas, which Plato has set forth in his most celebrated
dialogues, those that all agree in regarding as
the high point of his authorship: the \textit{Symposium},
the \textit{Phaedo}, and the \textit{Republic}. If these objections really
strike home, which seems to be conceded by Plato
himself, in that he leaves Socrates standing without any way out with respect
to his continued philosophizing, how then can it be explained
that Plato has set them forth? This question has
been answered in different ways. A whole series of German
scholars (\textsc{Schaarschmidt}, \textsc{Überweg}, and others) supposed that the \textit{Parmenides}
was not by Plato at all, because its content was incompatible
with the content of the generally recognized Platonic
dialogues. Others (\textsc{Zeller}, \textsc{Jackson}, \textsc{Natorp}, \textsc{Ræder})
sought to show that Plato himself, partly in the second section of the dialogue,
partly in later dialogues, has answered the objections put forward;
but in the content of the answer these scholars then differ
among themselves. Still others (\textsc{Apelt}, and recently \textsc{Wilamowitz-Moellendorff}
 in his biography of Plato) regard the objections
% --- p. 4 ---
\setcounter{footnote}{0}%
\opage{4}as practice examples, exercises for Plato’s disciples, Plato
having trusted that for the cleverest of these they would
fall away of themselves.

For my own part, after repeated occupation
with this remarkable work, I must agree with those who
think that the objections % print: Invendinger (misprint)
made in the first section really
strike home. But I think one must go a step further.
In these objections --- whether they derive
from others or from himself --- Plato has stood before points of view
that are of interest not merely for the criticism of his own philosophy,
but are of interest for every attempt at metaphysics
(in the sense of a conclusive knowledge of existence). I
shall seek to establish this through an analysis of the first
part of the dialogue. I remark at once that Plato has thereby grown
in my eyes. True, he has not been, and could not well
have been, conscious of the range of the objections put forward,
and they have not led him to give up the doctrine of Ideas entirely.
But within the remarkable mixture of criticism and mysticism
that constitutes Plato’s peculiarity, the critical element has in the \textit{Parmenides}
come forward at moments with such
great force that he later proved unable to bend it
into his system as a harmonious member, however significant
in several respects the new investigations may be that he
takes up in later dialogues. One understands how there could
issue from Plato’s school a critical thinker such as
Carneades, the Hume of antiquity, while Plato through the
other element in his philosophy became the progenitor of Neoplatonic
and medieval mysticism.

The second section of the dialogue will give a direction for continued
philosophizing, even if one concedes the objections to the doctrine of Ideas
put forward in the first section. The way that
is followed is that of hypothetical inference.
% --- p. 5 ---
\setcounter{footnote}{0}%
\opage{5}Different, mutually conflicting presuppositions are examined; their consequences
are to be decisive for which points of view one
adheres to. But the exposition is in most of its parts so
unclear, at the same time so abstract, and what is worse, so sophistical,
that people have found difficulty in ascribing it to Plato
or in regarding it as seriously meant. Those who reject
the authenticity of the dialogue draw mighty weapons from this
section. Those who already in the first section saw merely an exercise,
think that this holds to a still higher degree of
this second section. But there are also scholars who
have found here moments for a real answer to the objections of the
first section, in any case ways out
of getting around them.

For my own part I confess that I am unable
to find any significance in the greatest part of
the very abstract reasonings, in part built on decided fallacies,
that make up the second section of the dialogue.
But in the course of the unclear line of thought there appear remarkable
flashes of thought, which ought to be brought forward, and which in part
gain interest by being compared with other thoughts
in Plato. As a whole the section makes a very unfinished
impression. There is lacking both formal and real conclusion, and
it is not improbable that what we have here is really only
a draft or a scheme, which could have been rectified and made clear
by actual execution. Perhaps zealous disciples
have thoughtlessly published a draft left behind by Plato. ---

The conception of the two parts of this remarkable
dialogue thus indicated I now attempt to substantiate by
a closer review.

% --- p. 6 ---
\setcounter{footnote}{0}%
\opage{6}\chaphead{1. Critique of the Doctrine of Ideas.}


a. At a meeting in Athens between the Eleatics
Parmenides and Zeno on the one side, Socrates (as a
Platonist) on the other side, --- a meeting that
it is very difficult to make historically possible, --- there is read aloud
first \textsc{Zeno}’s famous (unfortunately now lost) work, in which
he seeks to demonstrate the contradictions in the assumption that
existence consists of diversities, and so is
a multiplicity. The famous objections to the possibility of
motion, which have especially attracted notice, are only consequences
of this criticism: motion presupposes multiplicity
in space and time, and the difficulty then is how
one can get across from one part of space and time to
another. After the reading Socrates says that Zeno’s
work contains the same as Parmenides’ didactic poem;
the difference is only that Parmenides straightforwardly declares that
only the One, that which is equal to itself and absolutely uniform,
exists, and that plurality and change are illusions,
whereas Zeno proceeds indirectly by showing that the assumption
of plurality and becoming leads to contradictions. This
Zeno concedes, and Socrates then maintains that the contradictions
that Parmenides’ opponents find in the doctrine of unity, and the
contradictions that Zeno finds in the doctrine of plurality, fall away
when one assumes that unity and plurality, likeness and difference,
rest and motion are so many Ideas, so many
% --- p. 7 ---
\setcounter{footnote}{0}%
\opage{7}kinds of being, and that the individual particulars (objects of experience)
each for itself in certain respects belong under
one of these kinds, in certain other respects under another,
(p. 127 E---129 E).

In this line of thought the opposition between
Eleaticism and Platonism appears clearly. According to Parmenides
existence was to be capable of being thought and to be thought through a single
thought, the thought of unity, of likeness (identity), and of immobility.
Plato’s doctrine of Ideas here marks an advance,
in that it maintains that existence is understood only through a plurality
of thoughts, whose mutual relation it is a matter of determining.
The doctrine of Ideas has come into being through the need to form
concepts for each individual delimited group of particulars, thus
through an ascent from a plurality of particulars to a unity
in thought, an ascent that had to be repeated every
time one stood before a new group of particulars. And Plato
now has Socrates maintain that each particular or each group of particulars
can fall under different thoughts or Ideas,
there being no contradiction in different Ideas meeting
in the same particular. The Eleatic critics do not go further
into this side of the matter, important as it is in and of itself.
What Plato here asserts can be described by what
H. C. Ørsted called multifoldness, by which he understood
the way in which different laws of nature in their interplay
make up the essence of the individual particular (phenomenon). The fountain,
e.g., becomes intelligible to us through the union of the laws
of fall, drop formation, light, and sound, etc. Ørsted
also uses the expression Idea of such a multifoldness.

The criticism that Parmenides now in what follows directs
against Platonism concerns, on the other hand, another main point,
namely the relation between Ideas and particulars in general.
Plato had maintained that existence cannot be thought merely
% --- p. 8 ---
\setcounter{footnote}{0}%
\opage{8}by means of the Eleatic One. Now Parmenides examines
whether existence can then be thought by means of the Platonic
Ideas.

First he asks whether Plato believes that there is any likeness
other than the one we find within our experience among
different particulars. To this Plato says yes. --- He has indeed,
we may interject, in the \textit{Phaedo} (p. 74) made use precisely of the concept of likeness
to show how the particulars of experience can present
different degrees of one and the same property, and how
by following the ascending series of degrees, here
thus of more and more complete likenesses, we can arrive
at the perfect, absolute concept, here thus the concept of
identity, as an ideal, for which experience can present only
steps or approximations.

From purely logical concepts, of which likeness is an example,
Parmenides passes over to ethical concepts, to the
beautiful, the good, and the just. Here too Plato says
boldly yes to the question whether there is something that is in itself
beautiful, good, etc., so that it is not merely experience
that shows us such properties. --- He had indeed also
in the \textit{Phaedo} (p. 75 C) expressly said that such concepts
as the beautiful, the good, the just must be regarded
in the same way as the concept of likeness. Experience
here too shows us only imperfect examples, but by
completing the scale we arrive at absolute Ideas.

Plato (the Platonic Socrates) stands differently when
Parmenides asks whether there are also Ideas corresponding to particulars
of a concrete or material nature, such as man, fire, water,
hair, dirt. Here he does not dare to answer yes at once, although he
sees that here too it is a matter of properties of given
particulars, properties that can appear in a more
or less perfect manner, so that the same line of thought
% --- p. 9 ---
\setcounter{footnote}{0}%
\opage{9}must find application here as with the logical and ethical
concepts. He stands in doubt on this point, which
Parmenides declares to be youthful uncertainty (p. 130
B---E).

This little passage has great interest in casting light
on the genesis of the doctrine of Ideas. More sharply than elsewhere (cf.
especially the passage cited in the \textit{Phaedo}) a distinction is here made between
different classes of Ideas. First the logical, then
the ethical fundamental concepts (categories) are set up, and last the
concepts of species that correspond to groups of individual particulars. We have
presumably hardly here an indication of the order in which the
different classes of Ideas dawned on Plato. For it
is probable (as also emerges from the dialogues
that are generally regarded as the first) that Plato
as Socrates’ disciple first caught sight of the significance of ethical Ideas
as guiding thoughts that had to have eternal and
unchangeable significance. Later the significance of the logical (epistemological)
Ideas became clear to him. In the \textit{Phaedo}
(p. 65 D) he accordingly begins with the ethical Ideas and
passes from them to the logical (the concept of likeness and identity),
whereupon it is said quite briefly that what holds
for such Ideas also holds for concepts such as magnitude,
health, strength, ``in one word, for the essence of all things.''
(Similar utterance p. 75 C). It is now some examples of
these ``all things'' with which Parmenides holds Plato’s feet to the fire.
The matter seems to be that Plato from the first did not
distinguish between such concepts as we form about the content
of our own intellectual functions (likeness, magnitude),
and such as make themselves felt in our ethical
evaluations, on the one side, and on the other side
such as serve to order and divide given
particulars (objects, individual beings), thus
% --- p. 10 ---
\setcounter{footnote}{0}%
\opage{10}concepts of species in a narrower sense. All his concepts he has won
by analysis and idealizing abstraction; but in the
dialogues that precede the \textit{Parmenides} he was not yet
attentive to the differing character of the concepts thus
won. Only later (in the \textit{Theaetetus} and in the
\textit{Sophist}) does he begin to single out a group of concepts
that correspond to what one may call fundamental and
formal categories, in that they are immediately tied to
the activity of thought itself and therefore find application everywhere
that we think, no matter what we think about.
He thereby laid the foundation of the doctrine of categories\footnote[1]{On Plato’s significance for the doctrine of categories see \textit{Den menneskelige Tanke,} p. 140 f.}. In contrast
to concepts that we apply in all thinking, and likewise
in contrast to those that we apply in the evaluation of actions
and persons, the proper concepts of species had then more and more to take their place, and we see accordingly that Plato in
dialogues from the following period (the \textit{Sophist} and the \textit{Statesman})
lays great weight on division and classification
of particulars.

In his thorough work, for me very instructive, ``Platon’s
philosophischer Entwickelungsgang,'' Dr. \textsc{Hans}
\textsc{Ræder} thinks that the doctrine of Ideas loses a part of its sublimity through
the concession Plato makes in acknowledging that even
dirt must have its Idea. The Idea of dirt, says Dr.
Ræder, cannot after all have been among the Ideas that
gods and men, according to the \textit{Phaedrus}, went out to behold.
To this it is to be said that the Idea of dirt is not itself dirty.
The concept or Idea of dirt would be the multifoldness of
elements and laws (chemical and mechanical laws, as we now
would call them) that conditions what becomes for us
% --- p. 11 ---
\setcounter{footnote}{0}%
\opage{11}the phenomenon of dirt. When Plato himself has felt misgivings
about this concept, it is because it is not a purely
intellectual concept, but at the same time a concept of evaluation: dirt
denotes something that one abhors, turns away from, or
gets out of the way. But it can very well be embraced with purely
intellectual interest, since one can disregard altogether the point of view of value.
This is nothing other than what Plato himself earlier
asserted. In the \textit{Republic} (p. 476 A) it is expressly declared
that the true philosophers form concepts not only
of justice, but also of injustice, not only
of the good, but also of the evil. Thus he has
himself here acknowledged an Idea of moral dirt. It
is then also from this side an inconsistency when he
for a time recoiled from acknowledging the Idea of dirt.
He ought consistently to have said with Heraclitus: ``Here too
are gods!''

b. That Plato has not been clear about the difference
between the different kinds of Ideas is due to the fact that
he has not distinguished between the different ways by which
he arrives at the different Ideas. There are two forms of
intellectual work that Plato employs without distinguishing definitely
between them.

The one consists in a going back, a regress from given
particulars to their presuppositions. It appears
clearly in the role that the concept of likeness, especially in its ideal
form as identity, plays with him. Parmenides had
already seen the fundamental significance of this concept; Parmenides’
``One'' means that which is in perfect identity with
itself, uniform in all respects and therefore continuous
(fr. 8. Diels.). This declaration has aptly been called
an act by which human reason took possession of itself
% --- p. 12 ---
\setcounter{footnote}{0}%
\opage{12}in that it asserts the supreme law of thought\footnote[1]{\textsc{Victor Brochard}: \textit{Études de Philosophie ancienne et de Philosophie moderne,} p. 133. --- Cf. also my treatise \textit{Parmenides fra Elea} (Mindre Arbejder. Anden Række) p. 141 f. (The contrast between what we now call formal and real science appears for the first time in principled form in Parmenides).}. It
was at once a protest against the multiplicity in the world of the gods
and against the chaos that human experience seemed to
present. It sets the sharp difference between the inner
coherence of thought, which in the end rests on the principle of identity,
and the manifold particulars of experience, which seem to
be capable of being connected in so many different ways. But
in Parmenides no distinction was yet made between unity,
identity, continuity, uniformity, and unchangeability; all
these expressions stood as designations for the same thing. Zeno
supplemented, as already touched upon, his master’s line of thought by
his polemic against plurality, discontinuity, and motion.
From the mathematical side Zeno is accordingly also now regarded
as the one who ``with his keen eye has pointed to the
difficulties that mathematics had to overcome in order to
become an exact science''\footnote[2]{\textsc{Zeuthen}: \textit{Hvorledes Mathematiken i Tiden fra Platon til Euklid blev rationel Videnskab} (Vidensk. Selsk. Skr. 1917) p. 15. See also \textit{Sur les livres arithmétiques d’Euclid} (Vid. Selsk. Overs. 1910), p. 432.}. It was by virtue of the ideal of identity and
continuity, dogmatically held fast, that Zeno
with so great energy stressed the impossibility of a rational
conception when plurality and discontinuity were presupposed.
Against this dogmatic conception Plato has the merit
of having maintained the necessity of more thoughts
than the one, if one wishes to understand the world. This is the significance
of his setting the world of Ideas in place of the Parmenidean
One. And quite a special merit he has in
investigating the ways by which we arrive at the thoughts
that make existence intelligible to us. Here belongs the
% --- p. 13 ---
\setcounter{footnote}{0}%
\opage{13}already mentioned line of thought in the \textit{Phaedo}, where it is shown that
the measure of likeness lies in approximation to ``the Idea
of Likeness,'' pure identity. Here belongs further
the line of thought at the end of the sixth book of the \textit{Republic}, where
he rises from the multiplicity and changeability of observations,
which at most make belief and conjecture possible,
to mathematical laws and from them again to the presuppositions
on which they, and with them all true science, rest, presuppositions
that it is the proper task of philosophy (of dialectic)
to bring forward in order to give them a conclusive grounding.
Through such analytic or regressive thinking
he reaches the stage on which Parmenides from the first
had placed himself, and which Zeno had sought to defend
through criticism of the opposite stage. But Platonism
and Eleaticism agree in the end in this, that the
ultimate ground for thinking is something that is, --- not
merely presuppositions and principles for our knowledge, but
the very own essence of existence. The world of Ideas is the highest
form of existence, just as Parmenides’ ``One'' was the
only form of existence.

The other form under which Plato’s intellectual work
appears is ordering and division. It is by this way that
he --- as one sees from the \textit{Parmenides}, after some hesitation
--- comes to conceive the Ideas as concepts of species,
each of which holds for a group of particulars.
This way he goes, e.g., in the fourth book of the \textit{Republic}, when
he distinguishes between different ``parts'' of the soul and builds
this division of the soul’s manifestations of life on their diverse,
even opposite or conflicting character.
In the \textit{Phaedrus} the definite demand is then made that one
shall be able both to draw together what is common to the
different particulars and to follow how this common element
% --- p. 14 ---
\setcounter{footnote}{0}%
\opage{14}appears under different forms. The division built upon this
is to be carried through until one halts at the
altogether simple, where no characterization is any longer possible.
In later dialogues (the \textit{Sophist} and the \textit{Statesman}) he describes
through examples from daily life and in broad and
copious elaboration the art of division, which both
lets likenesses and diversities come into their own. The
concepts of species formed in this way were for Plato,
like the logical and ethical concepts, ideal existences,
eternal and unchangeable beings. And just as all individual
likenesses in the world in the end found their explanation in
the eternal Idea of Likeness, so too the concepts of species contained
the explanation of the particulars for which they hold.
Thereby it is understood that Plato had no reason to distinguish
sharply between the different kinds of Ideas. What for
him was decisive was that they all contained the ultimate
grounds why existence in fact is as it is. If
we denote by the word ontology the line of thought according to which
the thoughts we form in our work of knowing stand immediately
and involuntarily as eternal and presuppositionless
beings, then Plato’s doctrine of Ideas is ontology just as much as
Parmenides’ ``One.''

c. We now return to the dialogue. --- The first objection
that Parmenides directs against the doctrine of Ideas, after he has
got Plato to conceive it in its whole consequence,
concerns the relation between Ideas and particulars. This relation
Plato had described as a participation in the Ideas on the
part of the particulars. All things that resemble one another are thereby
participants in the Idea of Likeness; all animals participants in the Idea of Animal,
etc. But, asks Parmenides, how can an Idea
subsist as the unity it is, when a multiplicity of
particulars is to partake of it? Each particular then gets only a
% --- p. 15 ---
\setcounter{footnote}{0}%
\opage{15}part of the whole Idea! But the Idea cannot be cut up into pieces, any more
than it can be present as a whole in any single
particular (p. 131).

Plato concedes the correctness of the objection. And that it
really strikes home lies in this, that just as general, eternal, and unchangeable
as the Idea is, just so concrete, perishable, and
changing are the particulars. This opposition can be overcome only
when it is seen that the Idea, the general concept, must
appear in a special, individual way in each single
particular. There is then no question of external division or of
community under the same cover. Justice, likeness,
man, fire, and dirt all appear in our experience
always in individualized forms. Aristotle was already
aware of this, and it was eagerly discussed in the scholasticism of the Middle Ages
and in modern philosophy (Spinoza, Leibniz,
and Berkeley). It is the problem of individuality that makes
itself felt. Plato’s great feat was to maintain the general
determinations of thought and their significance for our knowledge.
The next step, which could be taken only after long ages
of discussion, was to put such determinations of thought in the
right relation to the particulars for which they are to fit. Plato
rises with enthusiasm to his Ideas, but then afterwards does not
care how it goes with the particulars
he has used as starting points and steps in his
ascent. Nowhere, neither in the dialogue \textit{Parmenides} itself
nor in other dialogues, does Plato directly answer the objection
that has here come to expression, and which at a certain
stage of his intellectual life must have presented itself to him. But
his thoughts constantly circle around it, and it will become apparent
that it dawns on him that logical and mathematical thinking
can not only open to us access to the world of the pure
Ideas, but can also make it possible to find
% --- p. 16 ---
\setcounter{footnote}{0}%
\opage{16}thought-expressions for the changing forms of the particulars of experience, their waxing
and waning, their oscillations and their changes.
At the third objection we shall return to this
point.

Parmenides’ second objection amounts to this, that if
Ideas are to be necessary for the understanding of the particulars,
because these otherwise appear as an incoherent
multiplicity, then new Ideas are needed for the understanding
of the relation between those Ideas and the groups of particulars they
correspond to. In place of the incoherent multiplicity of particulars
there now meets us the incoherent multiplicity
of Ideas and particulars. The difference between Idea and particular must
pose the same problem as the difference between particular
and particular. We must, e.g., seek an Idea that explains the relation
between the Idea of Magnitude and the large things. There must
surely be something that is common to Magnitude and the large
things. That becomes an Idea in the second power. But with that
we are not done, for the question becomes how the relation
stands between the Idea in the second power on the one side,
the Idea in the first power and the particulars on the other side. We
thus enter into an infinite series and never become
done (p. 132 A, B).

This objection is often mentioned in Aristotle and has
led to a technical expression: ``the third man.'' Aristotle
takes for granted that it is known what is meant by this. Just as
man, as we know him in experience, presupposes
the Idea of Man (the concept of what a man is, which
holds for the individual specimens of man), so we must
by virtue of the doctrine of Ideas assume a new Idea of Man, which holds
for the relation between the first Idea of Man and
the specimens of man --- thus a third man, and
this can be continued to infinity. This objection is said
% --- p. 17 ---
\setcounter{footnote}{0}%
\opage{17}(according to a scholium of Alexander of Aphrodisias on Aristotle’s % print: Aristotetes’ (misprint)
\textit{Metaphysics}) to have been put forward by the sophist Polyxenus and
was probably known to Plato when he wrote the \textit{Parmenides}\footnote[1]{See \textsc{Baeumker}: \textit{Über den Sophisten Polyxenos.} (Rheinisches Museum. N. F. XXXIV).}.
Perhaps he has already in the \textit{Republic} X p. 597 C
tried to answer it. He says here that God has created
only a single Idea for each group of particulars, ``whether
he did not wish to produce more, or there was a necessity
to produce only one.'' Plato wants
to put an end to it, to shut off a line of thought that would
be able to continue itself to infinity; but he succeeds
only by a theological postulate. As later for Aristotle
(and still later for Hegel) the consequence that
a conception leads into an infinite series of thoughts is for him
a sufficient refutation of such a conception. It is
in reality the same line of thought that lies at the base
of Zeno’s criticism of plurality and motion.

For the theory of knowledge of modern times, by contrast, it stands as a natural thing that the work of knowledge cannot be concluded.
Even the last presuppositions to which the analysis
of our experience leads us back we dare not regard as
absolute and ``grounded in themselves''. There are series of thoughts
that can always be continued, and within which more intermediate terms
will always be found, however great a connection
has so far been found. The possibility of a series' constant
continuation, however many terms and intermediate terms are found,
involves for us no contradiction and no difficulty.
We hold to the law in accordance with which the constant continuation
is possible, and in this law we see the goal of knowledge.
And when it is a matter of presuppositions that are necessary
for gaining understanding, we see in the possibility of constant
% --- p. 18 ---
\setcounter{footnote}{0}%
\opage{18}continuation by virtue of the same law a proof that here we
stand at one of the peculiarities of the human faculty of thought.
The logical principles that are presupposed in
every attempt to ground an assumption cannot themselves be
grounded, since every attempt at it would presuppose them.
In a corresponding way, in the presupposition that
what holds for the magnitude n also holds for n + 1,
the fundamental peculiarity of mathematical
thinking\footnote[1]{Cf. on this \textit{Totalitet som Kategori.} (Vid. Selsk. Skr. 1917), p. 34---36.} has been found, particularly in its difference from formal logical thinking.
--- Already from the beginning of modern times great philosophers were
clear that in series that can always be continued
in accordance with the same law we must see the expression of a fundamental relation for
our knowledge, perhaps for existence itself. This appears
clearly in Bruno and Spinoza\footnote[2]{See \textit{Spinoza’s Ethica.} (Vid. Selsk. Skr. 1918), p. 8---9; 38 f.}.

\textsc{Wilamowitz-Moellendorff}, who on the whole does not attribute
great importance to the objections to the doctrine of Ideas, holds
that the second objection, the one I am discussing here, can easily be dismissed
by the following considerations. What several things have in common
is found by disregarding the differences between
these things and keeping to the points of likeness. The Idea
is thus formed by subtraction. But when one then again
compares the Idea with the things, there is nothing more
to subtract, thus no possibility of a new Idea\footnote[3]{\textit{Platon} II, p. 26.}. To this
it must be answered that precisely when there is nothing more to subtract,
we stand at just such a chaos of differences as it
was the task of the doctrine of Ideas to raise us above. A task
would thus in any case remain. But the main point
is that it is not in the spirit of Plato to regard the process of thought
that leads to Ideas as a mechanical isolating
% --- p. 19 ---
\setcounter{footnote}{0}%
\opage{19}procedure. The formation of Ideas is the fruit of a comparative
investigation that leads to a synthesis, or, as
Plato calls it, a synopsis, a seeing together, which leads the
manifold scattered phenomena back to one thought.
(\textit{Phaedrus} p. 265 C: εἰς μίαν ἰδέαν συνορῶντα ἄγειν τὰ πολλαχῇ
διεσπαρμένα). And one will not be able to get around the fact that a
new comparison is required with respect to the content between the
Idea thus formed and the particulars different from it, just as
a comparison was required with respect to the relation
among the mutually different particulars.

\textsc{Jackson}\footnote[1]{Journal of Philology X p. 298; 334.} is clear that Plato’s doctrine of Ideas, as
he has presented it in the \textit{Phaedo} and in the \textit{Republic},
is hit by the objection discussed here. But he lays
great weight on a way out that Plato himself for a moment tries out, % print: forforsøger (misprint)
namely, that one conceives the relation between
Ideas and particulars in such a way that the former are the eternal types, models
for the latter, so that what was called participation in
the Idea is in reality likeness to the Idea (p. 132 D).
Thereby one escapes, Jackson thinks, the mechanical
conception, according to which the Idea would as it were be parceled out
among the particulars. Jackson takes up more or less Plato’s theological
solution from the 10th book of the \textit{Republic}. When, in the production
of the particulars, certain eternal and unchangeable
types stood as guiding models, then this will make the ordering
and classification of the particulars possible, and the study of these will
lead us nearer to the knowledge of the Ideas as the perfect
expressions of what the particulars present only by approximation.


It is possible that this conception of the Ideas as models
was at bottom the one Plato in his innermost self favored.
In it, it seems, his philosophical and his poetic
% --- p. 20 ---
\setcounter{footnote}{0}%
\opage{20}need as well as his religious need could find satisfaction
at one and the same time. But Parmenides points out (p. 132 E) with
full right that the difficulty returns
quite in the same way as before, since the relation of likeness
between model and particular just as well must lead
to the assumption of an Idea as the relation of likeness between the particulars
among themselves. Here too we stand before the infinite
series, for which Greek thinking had such great
aversion. ---

Another way Plato has for a moment taken. There was
the possibility that the Ideas are only our thoughts, dwell only
in the souls, nowhere else. Then one would, it seems, escape
all disquieting consequences with regard to the relation of the
content of the Ideas to the particulars. But this possibility Parmenides
cuts off at once, and with full and quick assent from
Plato’s side, in that he maintains that every thought surely
must be a thought of something, and of something that has
being. A thought that is not a thought of something that is
(νόημα ὄντος) is a meaninglessness (132 B, C). --- It is a
highly characteristic statement, which throws light on
the doctrine of Ideas and on ancient thinking generally. One sees here
how near ontology lay to incipient thinking.
Thoughts may come and go, but the content of thought must
have an eternal and objective subsistence.

Now it is true that in the most recent times it has been wished to deny
that Plato’s doctrine of Ideas has this ontological character. It
has been asserted that the word ``being'' in Plato, when it is used
of the Ideas, means only the content of a statement\footnote[1]{``Nie und nirgends bedeutet `Sein' bei Platon, wenn von der Idee ausgesagt, etwas anderes als den Aussageinhalt''. \textsc{Natorp}: \textit{Platon.} (Grosse Denker. Herausg. von Aster. I.) p. 111. --- Already in his work \textit{Platon’s Ideenlehre} (1903) Natorp expressed this conception, to which J. A. \textsc{Steward} (\textit{Platon’s Doctrine of Ideas} p. 77) has adhered.}. And
% --- p. 21 ---
\setcounter{footnote}{0}%
\opage{21}on the basis of this assertion it has been held that whereas
the objection to the doctrine of Ideas that was called ``the third
man'' was irrefutable when the Ideas were conceived as
beings that stood in opposition to the particulars as
other beings, this objection falls away altogether
when one realizes that Plato’s Ideas are only methodological
points of view that lead to valid knowledge.

It is the viewpoint of modern theory of knowledge, founded by Kant,
that one wishes by this conception to read into
Plato. In a certain sense it is also not wrong
that the Ideas mean methodological points of view for Plato.
It is from the Ideas that he regards all particulars and passes
judgment on them, both as regards their reality and as regards their
value. It is --- as already indicated above ---
his significance in the history of thinking to have realized that
all our knowledge rests on presuppositions that it is a matter
of becoming conscious of. In his doctrine of Ideas he wished in his own
way to satisfy his master’s demand for self-knowledge.
But the point is precisely that the fact that the Ideas were methodological
points of view by no means excluded for him that they also
were beings that are. There was here for him no Either-Or,
but precisely a necessary Both-And. He did not distinguish
between validity and being. Therefore interest in the theory of knowledge and ontological dogmatism
could go together in him. He stands here at a similar standpoint with regard to his
Ideas as the Pythagoreans with regard to their numbers and the Eleatics
with regard to their ``One''. Modern mathematicians too
sometimes attribute ontological significance to the fundamental mathematical thoughts,
without its hindering them in the application of exact
mathematical methods.

For Kant there was here an Either-Or. It led him to
% --- p. 22 ---
\setcounter{footnote}{0}%
\opage{22}the declaration that the proud name ``ontology'' must be replaced
by the more modest ``Analytic of the Understanding''.

But even if it were justified to read the viewpoints of critical
philosophy into Plato, ``the third
man'' would not disappear. For the question must be raised
what the relation really is between our methodological
points of view and the particulars to which they are applied. Do
the particulars, with all their qualitative and temporal differences,
find their perfectly exhaustive expression in the
laws and equations to which the application of our methodological
points of view can lead? Is there here a relation of
identity or merely of analogy? Does there not in any case
always remain a problem here, of the same kind as the one that
had to go on existing for Plato in the relation between
Ideas and particulars? When Natorp denies the presence of this latter problem,
it is because he, as a representative of
the Marburg school\footnote[1]{Cf. my critique of the Marburg school in \textit{Totalitet som Kategori} p. 46---48. (Vidensk. Selsk. Skr. 1917).}, on the decisive point confuses
identity and analogy, just as Plato himself already
did.

For every world view --- theological, philosophical, or
scientific --- a problem arises here, which
always poses itself anew, and which arises from the fact that there is once
and for all a difference between observation (lived experience) and interpretation.
It may be that observation leads of necessity to a
certain interpretation; but as long as one cannot reverse the relation
and, starting from the thoughts the interpretation contains, derive
precisely the particular observations on which one built, so long
it cannot be proved that a relation of identity is present.
It is impossible to derive multiplicity from unity, whereas
analysis of the multiplicity may very well ground the
% --- p. 23 ---
\setcounter{footnote}{0}%
\opage{23}assumption of a unity (in the sense of lawlike connection).
In theological thinking the problem unity—
multiplicity appears in the relation between God and the world; but we
thereby get only the same problem again --- merely on a larger
scale\footnote[1]{Cf. my \textit{Religionsfilosofi} § 17, where Parmenides’ objection is already acknowledged and shown in its significance for theology.}. ---

If we return to Parmenides, he has his greatest
objection left, --- --- which will nevertheless prove to be closely
connected with the first two, to be as a continuation
of them. --- When the Ideas are as a world of their own,
with their own laws (αὐταὶ πρὸς αὐτὰς τὴν οὐσίαν ἔχουσι),
then they cannot be known by human beings, who belong
in the world of sensible experience, whose things are likewise
determined by their own mutual relations (αὐτὰ
πρὸς αὐτά ἐστιν). He who is master stands, e.g., in relation
to his slaves, not to the concept (the Idea) slave, and conversely,
whereas the concept (the Idea) master and the concept (the Idea)
slave stand in relation to each other. And so the concept
(the Idea) science relates to the concept truth (the Idea
of truth)\footnote[2]{Cf. with this \textit{Phaedrus} p. 247 E (In pre-existence the souls behold absolute science: --- τὴν ἐν τῷ ὅ ἐστιν ὄν ὄντως ἐπιστήμην οὖσαν).}, but our science relates to the
truth we can attain. The Ideas do not dwell in us, after all, but in their
own world, and the world of Ideas cannot make intelligible to us a world
different from it. Not even God, who after all
must be presupposed to stand in relation to the Idea of science\footnote[3]{That God is God is due to his relation to the Ideas. Cf. \textit{Phaedrus} p. 249 C (What the souls saw in pre-existence, the Ideas, is that, πρὸς οἷσπερ θεὸς ὤν θεῖος ἐστιν).},
would be able to know the world we live in, and just as little
will we from our world be able to ascend to the world of Ideas.
% --- p. 24 ---
\setcounter{footnote}{0}%
\opage{24}Neither being an image nor participation can here
come into question (p. 133 B—134 B).

That this objection hits the mark can be seen from the theory of knowledge
Plato himself developed in the sixth book of the \textit{Republic}.
Within the world of thought (τὸ νοητόν) a distinction is made between
mathematics, which works with certain presuppositions whose
grounds are not given, and dialectic, which uses
hypotheses only as steps and springboards to reach the unpresupposed,
the principle of everything, and this highest knowledge proceeds
from Idea to Idea and ends in Ideas (\textit{Republic} VI p. 511 B).
The doctrine of Ideas is thus a closed whole, which does not point beyond
itself. Once one has merely found the entrance to it, one goes
forward along the path of pure deduction, keeping to what
follows from the principle (ἐχόμενος τῶν ἐκείνης [τῆς ἀρχῆς]
ἐχομένων), and one ends in Ideas (τελευτᾲ εἰς εἴδη) and
does not get out of the world of Ideas. --- It stands here as
incomprehensible how this pure deduction can acquire any
significance for the world of experience, which Plato nevertheless expressly
presupposes, since he lays weight on the fact that only those
who have attained this highest knowledge are suited to
lead human affairs. But he has not shown
any possibility that pure deduction could come
into any interaction with the knowledge that is won by
observation.

Parmenides’ third objection now draws precisely only the
consequence of the doctrine of Ideas as purely deductive knowledge.
But he goes a step further, in that he denies the possibility
of reaching the point where this deduction could begin.
Not even the steps and springboards mentioned by Plato become possible,
since the world of experience too is closed in itself.
Where within the world of experience would that point lie
that could make an ascent, a leap up into the world of Ideas
% --- p. 25 ---
\setcounter{footnote}{0}%
\opage{25}possible? --- The line of thought is the same as the one Kierkegaard
applied against Hegel\footnote[1]{Cf. ``\textit{Søren Kierkegaard som Filosof}''\textsuperscript{2} p. 61 f.}, in that he demanded that ``the
dialectic of the beginning'' be made clear before one began
on the system. How is the transition made from analysis of
experience to the deductive synthesis? Kierkegaard draws from this
the conclusion that there is no presuppositionless philosophy. ---

Already the first two objections touched essential
sides of the doctrine of Ideas. But the third objection draws
the ground clean away from under it. It redounds to Plato’s honor
that he has let the radical misgiving about the doctrine of Ideas
come so clearly to expression. One may say that he has here
come right to the point where Kant in the year 1772 found
the foundation of his decisive formulation of the problem. In his letter to
Marcus Hertz of 21 Feb. 1772 Kant asks precisely how
it can be that the concepts we form by our mental
activity can acquire significance for a world we have not
ourselves produced\footnote[2]{See \textit{Kontinuiteten i Kant’s filosofiske Udviklingsgang} (Vid. Selsk. Skr. 1893) p. 52.}. In Plato the question is formed thus:
how can human science become divine
science? ---

Just as little as the first objections does this
third objection lead Plato to give up the doctrine of Ideas, and we shall
presently see why he does not give it up.
But this third objection does lead him to make
changes in his philosophy. It is the relation between deduction
and induction that is decisive. These two methods
go their separate ways in Plato in the great dialogues (\textit{Symposium}, \textit{Phaedo}
and \textit{Republic}). The task had to be to find
a possibility for a cooperation between them. On such
a cooperation modern science builds from the
% --- p. 26 ---
\setcounter{footnote}{0}%
\opage{26}Renaissance onward, and as we shall see, Plato’s
later thinking too moves in this direction, even if
it does not reach the goal. ---

After the three objections have been discussed, Parmenides says
that such difficulties are what one comes to when
one attributes being to the Ideas. But if on account
of these objections one wishes to give up the doctrine of Ideas, then one will
be unable to give one’s thinking any definite direction
(οὐδὲ ὅποι τρέπει τὴν διάνοιαν ἕξει) and the faculty of intellectual
discussion (ἥ τοῦ διαλέγεσθαι δύναμις) will in the end
disappear (p. 135 ABC). --- More strongly and more characteristically
it could not be expressed that ontology had to
be held fast, and why it had to be held fast. Parmenides
and Plato are here fully agreed. And Parmenides is in general
in this dialogue Plato himself, or a part of him. One
can speak of the Platonic Parmenides just as of the
Platonic Socrates.

\chaphead{2. The Second Part of the Dialogue.}


When the young friend of the Platonic Parmenides (the Platonic
Socrates) is in perplexity as to what is now to
become of philosophy, the old philosopher first acknowledges
his interest in the work of thought, in contrast to those who
let thought roam about without law or rule. And thereupon
he urges him to exercise his thought by starting from
sharply determined presuppositions. In order to find where the problems
really lie, one must investigate both what
follows from a certain determinate presupposition, and what
follows from the opposite one. Now unity and plurality, likeness
and unlikeness, rest and motion, coming-to-be and dissolution,
being and non-being stand opposed to one another as opposites.
We must investigate what the consequence will be of holding
% --- p. 27 ---
\setcounter{footnote}{0}%
\opage{27}to a single term within each of these pairs of opposites.
Only by such an all-sided discussion can one attain
truth and acquire reason (135 C—136 E). Parmenides
will then in what follows discuss his own presupposition of
the One as presupposition for everything (137 B).

The treatment of the problem proposed by Parmenides has
an evident occasion in the preceding investigations. The difficulties
of the doctrine of Ideas proved everywhere to rest on the opposition
between Idea and particulars. The Eleatics escaped this
problem in that they established in advance that only ``the One''
(corresponding to Plato’s world of Ideas) has reality. It is a
great concession Plato lets Parmenides make when he
lets him go into discussing both the presupposition
that ``the One'' is, and the one that it is not. Yet this
second presupposition, that ``the One'' is not, that thus
the manifold is, had already been discussed by Zeno in his
famous work, whose line of thought is also traced in many
places in what follows. But Plato lets Parmenides introduce
new distinctions, which conflict with the Eleatic presuppositions
and make the discussion very involved, not to
say wholly obscure. First of all we must say that the subject of the
proposed method was not happily chosen. It was after all Platonism,
not Eleaticism, that was to be examined. True,
Plato agreed with Parmenides that the valid content
of our thoughts must mean something that is, and also that thought
seeks to connect its manifold content as intimately as
possible, thus to find unity and coherence in it. But
Plato’s advance over Parmenides consisted precisely in
Plato’s with great energy and enthusiasm both demanding and
describing the spiritual work that could lead to finding
this unity, which Parmenides simply presupposed. In
the \textit{Meno}, \textit{Symposium}, \textit{Phaedo}, \textit{Republic} and \textit{Phaedrus}
% --- p. 28 ---
\setcounter{footnote}{0}%
\opage{28}he had depicted the development from the sensible absorption
in the stream of particulars to the unity and inner coherence of the life of thought.
And it had been shown that ``unity'' means lawlike
coherence, rationality or identity. In the \textit{Parmenides}
p. 139 D it is expressly said that unity and identity are
the same (ἥπερ τοῦ ἑνὸς φύσις, αὕτη δήπου καὶ τοῦ ταὐτοῦ),
and with this statements in the \textit{Republic} (p. 533 C; 537 C)
and in the \textit{Phaedrus} (p. 265 DE; 266 B) agree. It would have been
both more natural and more fruitful if it had been
the concept of identity, and as its opposite the concept of difference,
that had become the object of discussion, so that
the consequence of laying only the one or the other of
these two opposed terms as ground could be brought out.
The problem unity-multiplicity is derivative in relation to
the problem identity and difference. What Zeno was after
in his famous aporias was, as already mentioned, really
not motion, but plurality, absolute multiplicity
(what I have called the chaotic series of differences), which
excludes understanding. It is --- as Plato has clearly seen ---
the possibility of understanding that is the Eleatics’ real problem,
which Parmenides discusses positively, Zeno negatively. For Plato
it would have been most natural to let the discussion turn
on the possibility of understanding itself, thus of science.
Because the concept of unity is pushed forward before the concept of identity,
it becomes more Eleaticism than Platonism that is
discussed in what follows, and thereby the discussion becomes both
unfruitful and unclear. It is at bottom neither Eleaticism
nor Platonism, in general no clear position, that is examined.


After often repeated occupation with the second part of the dialogue
I have for my part gained a decided impression of
a lack of natural design, and at the same time of defective
% --- p. 29 ---
\setcounter{footnote}{0}%
\opage{29}execution in the particular and of a lack of natural conclusion.
It had after all been announced that a thorough thinking-through of
the various possible hypotheses was to help us get going
again in our philosophy, when the hope we had cherished for the doctrine of Ideas
in its hitherto existing form had to be given up. But there
is no indication whatever in this respect.

As already mentioned, here lies one of the reasons why
many for a time were inclined to declare the dialogue
spurious. The obscurities and fallacies that the second part presented
one would not put to Plato’s account. Other scholars
have conceived the second part as a youthful work
composed for polemical purposes, teeming with sophisms; thus
Otto Apelt (Beiträge zur Geschichte der griechischen Philosophie
p. 59; 66). \textsc{Wilamowitz} in his biography of Plato
agrees in part with Apelt. He regards the dialogue as unfinished,
produced in a period of Plato’s life when, after
practical and theoretical disappointments, he tried new paths. For my
part I too consider it probable that the final touch
has not been put on the second part of the dialogue, which
contrasts so strongly with the carefully elaborated
and thoughtful first part. I therefore surmise that Plato on
closer reflection was led to pose the main problem
that now presented itself to him after the severe self-criticism,
in a way that was both clearer and better corresponded
to his presuppositions. The following dialogues have, as
I shall show in a later section, decided indications in the direction
of such an altered formulation of the problem.

With regard to the course of the second part of the dialogue I will
refer to the scholars who believe they have found a positive
content in it, such as \textsc{Eduard} \textsc{Zeller} (Platonische Studien
p. 174—179), \textsc{Jackson} (Journal of Philology XI p. 306—
331) and \textsc{Ræder} (Platon’s philosophische Entwickelung
% --- p. 30 ---
\setcounter{footnote}{0}%
\opage{30}p. 309—316). These authors have found as the general result
of this second part of the dialogue the knowledge
that unity and multiplicity constantly presuppose each other,
so that one cannot exclusively lay one of the two opposed terms
as ground for one’s thinking. It is only one of the
many examples of the fact that our concepts stand in relations of correlation
and therefore serve for mutual illumination and determination
of one another. But this truth Plato would certainly,
had he worked through his draft, have been able to
develop in a clearer and more natural way and in
closer connection with his fundamental thoughts.

Only a few individual points shall I bring forward here as
a characterization of the second part of the dialogue. It is partly a sharply
carried-through thought, partly a wondrous sophism, partly a
significant flash of thought.

Parmenides begins by examining his own fundamental thought,
the thought of ``the One''. If the concept ``the One'' is
understood unconditionally, no predicates at all can be attributed to
it. It cannot be called whole, since it can have no parts,
and when it is no whole, it can have neither beginning
nor end, thus cannot be in time and cannot undergo
any change. It cannot have any location,
since something outside it would then be presupposed. It cannot
be in itself, since it would then be both greater and
smaller than itself, nor be in anything else, since
there would then be something else besides it. It can just as
little be at rest (``stand'') as be moved, just as little
be identical with itself as different from itself.
But neither can it be identical with anything else
or different from anything else. (It is in this connection
that it is betrayed that unity and identity are the same). And
if the predicate being means the applicability of one of the
% --- p. 31 ---
\setcounter{footnote}{0}%
\opage{31}concepts mentioned, in particular of temporal determination, then one cannot
call it that which is. No predicate can thus be applied
to it, and therefore it can neither be named, sensed nor
thought (137 C—142 A).

Plato here has Parmenides draw the ground away from under
his own doctrine. In his didactic poem the historical Parmenides
--- as already stated above --- had attributed to his ``One''
several predicates: it is One with being; it stands in relations of likeness
and continuity to itself; it has circular form.
The paradoxical thing is that the concept of unity (identity) was precisely
set up by the Eleatics as the principle of the intelligibility of existence.
And now it turns out that it is itself unintelligible!
--- When this had been shown, there might have been reason
to examine the concept of intelligibility itself. That is the
way by which Plato has distanced himself from Eleaticism. And
it led him to a plurality of Ideas, corresponding to different
groups of particulars. In the first part of the dialogue it was after all
expressly maintained that there is no contradiction in the fact that the same
particular can belong under different Ideas, thus have
different predicates. This line of thought Plato could
have carried further (and he does so in later dialogues); but
it is not alluded to at all in the second part of the dialogue.

With regard to the arrangement within the dialogue one must
wonder that the critique of ``the One'' has not found
its place in the first part, as a link in the first objection to
the doctrine of Ideas. For the critique of the ``One'' acquires significance
not only with respect to Eleaticism, but also with respect to Platonism,
since each single Idea stands as a unity in opposition to the
multiplicity of particulars that belongs under it. Moreover
Plato (as the \textit{Republic} p. 507 and the Seventh Letter
show) never gave up the hope of being able to lead all Ideas
% --- p. 32 ---
\setcounter{footnote}{0}%
\opage{32}back to a single fundamental Idea, and here the problem
unity-multiplicity comes forth again.

But not only with respect to Eleatics and Platonists does this
remarkable critique of the concept ``the One'' have its significance.
Again and again, both in the history of religion and in that of thinking,
attempts appear to gather everything in one thought or
to start from a single thought. All mysticism and romantic
speculation is hit to that extent by this critique. Yet
consistent mystics will admit its result, in that
they maintain that if God is the absolutely One, nothing can be
said about him, indeed, he is nothing, i.e. nothing of what one
can name or think\footnote[1]{Cf. \textit{Religionsfilosofi.} § 21. Conclusion.}.

Plato himself has paved the way for the right understanding
of the mode of operation and task of human thinking, in that
he (in the \textit{Republic} and in the \textit{Phaedrus}, and later in the \textit{Theaetetus}
and in the \textit{Sophist}) teaches that all thinking is a comprehending
through the interaction of connection and separation. He is
thereby a predecessor of the modern theory of synthesis in
Leibniz, Kant and their successors. ---

The next consideration (p. 142 B--155 E) proceeds from a
presupposition that really ought to have been rejected, if one
trusted the result of the first consideration.  For, as
we saw, the predicate ``being'' (οὐσία) had to be rejected if
``the One'' was to be taken seriously.  But the new consideration
proceeds precisely from the premise that the One ``is.''  The first
consideration proceeded from the premise that what is, is ``one'' (εἰ ἕν ἐστιν);
the new consideration proceeds from the premise that only ``being''
can be attributed to the One (ἕν εἰ ἐστιν).  There, the starting point was
the content of the concept; here, the starting point is that this content
has existence, --- that ἕν is an ὄν.  The discussion of this presupposition
leads to results opposite to those of the discussion of the
% --- p. 33 ---
\setcounter{footnote}{0}%
\opage{33}former.  That the presupposition is nevertheless taken up and gone through at great length is evidence of the provisional
and tentative character of this second part of the dialogue.  Moreover, the discussion
is also, to be sure, evidence of what sort of sophistry
Plato could fall into before his thought had attained
full clarity.

The line of thought is this.  When the One is said to be,
we get two concepts, the One and its being.  But these
concepts must belong closely together and make up a whole.
Every whole consists of parts, and each of the parts
must in turn have the character both of unity and of being,
and these two predicates must in turn form a whole --- and so
on to infinity by virtue of the line of thought once embarked upon
(κατὰ τὸν αὐτὸν λόγον ὡσαύτως).  But it is
not merely a duality that thus repeats itself; it is a
triad, for the two predicates and their connection make up
a triad, --- and here too we can go on in the same
way, so that we come to need the whole series of numbers.  The One
then becomes a multiplicity.  And as a whole it is complete,
has both beginning and end.  It is, moreover,
both at rest and in motion, advancing through different
moments, at each point standing between a
``now'' and a ``thereafter.''  --- Since there is thus the possibility of a
whole series of predicates, each of them separately demonstrable, observation
and science can now become possible, which they
could not after the former consideration. ---

Even Zeller and Jackson, who otherwise know how to find deep
philosophy in this section of the dialogue, give it up here.
There is also an obvious sophism in the manner in which
the concept of a whole is introduced here.  The fact that a distinction is made between
two concepts or predicates (unity and being) shows
in a way that there is a whole, since predicates that
% --- p. 34 ---
\setcounter{footnote}{0}%
\opage{34}are attributed to the same particular must stand in a certain connection
with one another.  In the first part of the dialogue Plato has indeed also
maintained that there is no contradiction in the same particular
falling under several Ideas.  It is admittedly an indeterminate
whole; it is the one that appears in the third Aristotelian
figure of inference (B is A and B is C), where the difficulty
is precisely that of finding what connection one
has a right to assume between different predicates of the
same subject (between A and C).  It is in any case
a purely logical whole, a connection of thought, and it is
unwarranted to transfer to it what holds for wholes
in time and space, wholes that are given in actual
or possible experience, --- real wholes.  And yet it is
assumed without further ado that the whole that the concepts unity and
being make up must be subject to the categories of space, time, and number.
From the standpoint of the historical Parmenides, according to which
the ``One'' is conceived from the outset as a uniform
sphere throughout, the category of space is in any case applicable;
but its applicability is not reached by inference, but is
presupposed.  To Plato this line of thought is wholly foreign,
apart from the confusion of concepts by which it is justified. ---

It is no accident that the line of thought proceeding from the ``One''
as being (existing) takes on a realistic character.
When we ask what we mean by saying that something is, exists,
psychological investigation already shows that we thereby
rely on observations which prove that this something stands
in firm and unbreakable connection with everything else that is present
to us.  If any of this other is doubted, the
same question is simply raised anew.  To be (to exist) is
to be a member of a lawful connection.  When it
is therefore to be discussed what lies in the ``One's'' being,
one must necessarily become clear that one thereby in advance
% --- p. 35 ---
\setcounter{footnote}{0}%
\opage{35}places it in relation to something else, and this is understood
to the extent that the line of thought slides over into the application of
concepts such as space, time, and number, which are precisely forms of
connection.  One enters the world of relations
as soon as one asks about existence.  And even if one confines
oneself to the mere content of the concept, as in the purely formal
sciences, one does not avoid the relation to something
else.  Logically and mathematically everything is possible that does not
contain self-contradiction; but even this formal criterion of truth
presupposes that the ``One'' can be put in relation
to something else; otherwise its logical validity
cannot be established.

One might ask whether the view put forward here concerning ``being''
and existence is not foreign to Plato.  To this the answer must be that Plato himself (in
the tenth book of the \textit{Republic}) applies the same criterion
to being as is put forward here.  It is on the occasion
of the relation between art and science that Plato comes
to this point.  Art keeps to the single
image as it appears when the object is seen from
a single, definite side.  Science arrives at the true
concept of the object by relying on observations of
the object from as many different sides as possible.
One and the same table looks different to us, according to
the side from which, and the distance at which, we see it.
A stick that under ordinary conditions is seen as straight will,
when it stands in water, be seen as bent.  Under different
conditions one and the same thing can appear as soft and as
hard.  If we want to know the real nature of things,
we must count, measure, and weigh; we can then form a concept
of them, on the basis of which, taking the
different circumstances into account, we can explain how it
% --- p. 36 ---
\setcounter{footnote}{0}%
\opage{36}comes to look different.  Being we attribute to what
a concept thus formed contains.  (\textit{Republic} X p. 598 A;
602 C--603 A.  Cf. VII p. 528 AB).  --- We shall later have occasion
to see that Plato in his later writings places great
weight on this consideration, which echoes in Galileo's
famous exhortation to measure everything that is measurable, and
to make measurable what is not yet so.  It is very
remarkable that Plato has not made use of it in the \textit{Parmenides}.
It would have supported the critique of the concept of
the ``One'' that he has given beforehand, and would have led to the
result that later works its way forward in the \textit{Philebus} and in the
\textit{Timaeus}, namely, that the mystical concept of the ``One''
acquires rational significance only by being seen as an expression of
the formal and real connection of the particulars.  Spinoza's ``substance''
is precisely Parmenides' ``One'' as the principle of the world's
inner, lawful connection\footnote[1]{\textit{Spinoza's Ethica} p. 9--11; 14--18. (Vidensk. Selsk. Skrifter 1918).}. ---

The last section of the dialogue (p. 164 B--166 C) is a counter-test
of the first consideration (which proceeded from the validity of the One).  It proceeds from the premise that the ``One'' is not, and
investigates what the consequence of this will be for the many
(τὰ πολλά).  The many particulars then stand merely in a relation of difference
to one another, and there will be absolutely no connection
between them.  Even if they apparently form
quantities or masses (ὄγκοι), on closer consideration these will
dissolve into chaos, into an infinity of different
particulars.  And in every apparent whole one will find
ever new parts within it and ever more extremities at
its boundary.  In the end multiplicity itself will vanish,
since its parts would each of them have to be ones, but no
One exists! ---

% --- p. 37 ---
\setcounter{footnote}{0}%
\opage{37}It can be clearly traced here that the meaning of ``the One''
is that it is the principle of lawful connection in the world.
If this concept falls away entirely, we arrive at
a limit of our knowledge that is the opposite of the one to which
the first consideration (137 A--142 A), which discussed the content
of ``the One,'' led.  Just as impossible as it is to
hold fast to the concept ``the One'' in its full consequence,
so impossible is it to give it up entirely.  There are thus two
mutually opposed presuppositions, each of which taken alone would
make knowledge impossible.  In my presentation of the theory of knowledge I
have called the one ``the absolute series of identity,''
the other ``the chaotic series of differences''\footnote[1]{\textit{Den menneskelige Tanke} p. 162--172. --- See also \textit{Totalitet som Kategori} p. 8--11. (Vidensk. Selsk. Skr. 1917).}.
Between the two all human knowledge moves, in
constant conflict of identity and difference. ---

The dialogue ends suddenly after the discussion last mentioned,
without any retrospect, without concluding remarks,
without any hint whatever of the interest that the
various presuppositions discussed might have for the question dealt with in the
first part of the dialogue, that is, without
pointing to what might take the place of the so severely
criticized doctrine of Ideas.  This speaks strongly in favor of the supposition
that we have here only a loose sketch. ---

I have here discussed only the beginning of the second part (p. 137 C
--155 E) and its end (p. 164 B--166 C).  The considerations
that lie between (p. 156 A--164 A) I do not go into (apart from a
single point).  I cannot find anything
fruitful in them, and the confusions of concepts flourish more luxuriantly
there than in other parts of the dialogue. --- The single
point I wish to dwell on requires a section of its own.
% --- p. 38 ---
\setcounter{footnote}{0}%
\opage{38}\chaphead{3. The Concept of the Sudden (τὸ ἐξαίφνης) in Plato's Philosophy.}

The investigation of ``the One'' as being (142 B--155 E)
had led to a closer development of a thought that
the Platonic Socrates had maintained right at the beginning
(129 B), namely that several Ideas may very well at one and the same time
hold of the same particular, a piece of knowledge which, moreover,
had already been won in earlier dialogues (\textit{Phaedo} 102 C;
103 B. --- \textit{Republic} 479 A; 523 C).  And in a remarkable passage
that concludes the investigation mentioned, Plato
(the Platonic Parmenides) now comes to the question of
a successive relation to the Ideas that may hold of the same
particular, --- of the transition from likeness to unlikeness, from
rest to motion, from smallness to greatness, etc.  And
here the concept of the sudden emerges, which presents
a special interest in Plato's philosophy.

a. The introduction of this concept itself comes very
suddenly.  In the demonstration that ``the One'' as
being must be able to have many different predicates,
it is particularly emphasized that it must be subject to time;
otherwise it could not pass from one state to another.
But emphasis is laid --- with a certain concession to Zeno
--- on the point that in every now it must be what it is, the now
cannot be passed over, notwithstanding that one must constantly advance from
moment to moment, so that the now must at once be let go again and
the next now (τὸ ἔπειτα) seized, (152 CDE).  Here no serious
problem is as yet hinted at.  But then, a little later
(155 E), the question is taken up again, in that the opposition or
contradiction between the various predicates (or the states
these indicate) is sharply emphasized.  A change (μεταβολή) must take place, a transition from one state to a
wholly different state.  The one state must be let go, the
% --- p. 39 ---
\setcounter{footnote}{0}%
\opage{39}other taken over.  But when does this transition take place?  As long as
the earlier state (likeness, rest, smallness, etc.)
is there, it does not take place, and when the new state (difference,
motion, greatness, etc.) is there, the transition has already
taken place.  It takes place suddenly.  But this sudden (τὸ ἐξαίφνης)
is a paradox (ἄτοπον), for it means a transition between
two states, and since no time can be given in
which a particular is neither at rest (etc.) nor in motion
(etc.), the sudden is in no time.
Into this sudden the old state passes, and out of
this the new state emerges.

The solution that is hinted at here consists in acknowledging
the reality of the transition, and with it of the sudden, but in
placing it outside time.  What is outside time and
yet has validity must, according to Plato's philosophy, belong
to the world of Ideas, and Plato would consistently have had to assume
an Idea of suddenness.  He would surely have been even more reluctant to do so
than to assume an Idea of dirt.  For this
new Idea would bring unrest into the calm world of the Ideas.  In
a double sense the sudden then becomes something that has
no home anywhere (an ἄτοπον); neither the world of the Ideas nor that of the particulars
can house it.  But on closer examination
Plato's thinking does not entitle him to place the
sudden outside time.  Though it is not supposed
to be a state, it nevertheless fills a link in the series of states, and is
thus a particular (a phenomenon, something that is present).  The ground
for its timelessness rested only on the transition's being
neither rest (etc.) nor motion
(etc.).  But the error of the Platonic Parmenides might
precisely lie in his presupposing a contradictory relation
between these concepts.  When the same particular can well
fall under different Ideas, a Both--And is not
% --- p. 40 ---
\setcounter{footnote}{0}%
\opage{40}excluded in advance; likeness and difference, rest and motion,
etc., need not be mutually contradictory
concepts.  It is not necessary that one Idea disappear
in order that another may come to hold.  The transition can
thus take place from a state of likeness with subordinate
diversities to a state of diversity with subordinate
likeness, --- from a restless rest to a uniform motion,
etc.  In more recent times Descartes still proceeded from a
contradictory content between rest and motion, but
Leibniz showed the possibility of a transition from rest to motion,
when rest was conceived as a state of equilibrium,
in which different tendencies of motion held one another
in check.  But just as the question of ``being'' (existence)
could not be answered by keeping to a single being
taken by itself alone, but only by seeing it in relation to other
beings, so the transition from one state to
another cannot be understood from the single state as isolated;
only under certain definite conditions, in particular under the
influence of releasing impulses coming from outside, does
the transition take place.  It is isolation that must be abolished,
within the world of Ideas as well as within that of particulars. ---

In two modern thinkers Plato's doctrine of the
sudden has aroused interest.

Kierkegaard, who in a note to \textit{The Concept of Anxiety} (ch. 3)
discusses this doctrine, translates ``the sudden'' as ``the moment''
and says that for Plato it is ``the non-being
under the determination of time.''  This Plato would surely
not concede.  He knows no empty time.  He
expressly maintains, after all, that the sudden is not in any
time.  But Kierkegaard is right to the extent that Plato
cannot consistently place the sudden outside time, since that
to which a state passes over, and from which another
% --- p. 41 ---
\setcounter{footnote}{0}%
\opage{41}state begins, must itself be a link in the series of time.  When
Kierkegaard without further ado translates (τὸ ἐξαίφνης) as ``the moment,''
he has thereby already corrected Plato; for
every moment is a link in the series of time.  Moreover,
Kierkegaard needed a positive meaning and a definite
temporal location for the ``sudden'' in his doctrine of ``the leap.''
The leap, the absolute act of freedom, is to be a positive
action at a definite moment and is to be capable of being observed.
But it is not, like other states or moments,
to be determined by law; it is not subject to certain definite
conditions of release.  ``The sudden,'' says Kierkegaard
(\textit{The Concept of Anxiety}, p. 137), ``knows no law.''  As I have
shown in my book on Kierkegaard as a Philosopher (second
edition p. 78--82), Kierkegaard nevertheless locates ``the leap''
between two moments; in the first it has not yet
taken place, in the second it has already taken place.  Hereby he
again approaches Plato's own conception of the sudden.  But
then the sudden cannot be called ``the moment,'' if by
a moment one understands an independent link in the series of time.
The two moments between which the leap, according to Kierkegaard,
takes place have no moment between them.  Strangely
enough, Kierkegaard holds that the leap is nevertheless to be capable of being
observed.

Natorp has seen that according to Plato the sudden is not
in time, and he finds in this a support for his interpretation
of Plato's doctrine of Ideas in analogy with modern theory of knowledge
(especially that of the Marburg school).  He finds in it, namely,
the conception expressed that not sensation, which is bound
to time as discontinuous succession, but only thought
can maintain continuity, can ground the transition from one predicate
to another.  Whereas Kierkegaard interprets the sudden
as the leap, Natorp without further ado sets ``the transition's
% --- p. 42 ---
\setcounter{footnote}{0}%
\opage{42}continuity'' as synonymous with ``transition''; that the transition
takes place ``im Auf-einmal'' (Parm. 156 D) is supposed to mean
that ``the continuity of the transition has its root in a thinking''\footnote[1]{The section ``Platon'' in Aster's ``Grosse Denker'' I, p. 127.}.
But Plato quite evidently regards the sudden transition
as a fact that it is a matter of understanding, and his ground
for its not being able to be in any time is purely negative, without
any reference to thought, to the world of Ideas.  It would
likewise be of no use to refer the transition to pure logic,
partly because there it becomes continuous transition, whereas
the riddle is precisely the discontinuous in the transition from Idea
to Idea, partly because a new problem then emerges,
namely what the relation is between the logical continuity
(which is really identity) and the real (temporal)
discontinuity.  But this is, to be sure, precisely a question
whose significance the Marburg school overlooks\footnote[2]{Cf. \textit{Totalitet som Kategori} p. 46--48. (Vid. Selsk. Skr. 1917).}.  This
school here faces the same difficulty as Parmenides
and Plato, an absolute opposition between continuity
and discontinuity, and it underestimates the constant
significance of discontinuity in our knowledge as a spur
and as a factor that poses problems.

The fact of discontinuity, the suddenness with which
a transition can take place, is closely connected with the fact that
every particular stands in relation to other particulars.  If we
had at our disposal only a single experience, the concept of
discontinuity would not have been formed.  A single point taken by itself
alone cannot be called discontinuous.  And suddenness cannot
be predicated of any experience that does not stand in qualitative
opposition to other experiences.  A closer investigation
shows that our knowledge develops under constantly renewed
conflict between continuity and discontinuity.  To this
% --- p. 43 ---
\setcounter{footnote}{0}%
\opage{43}conviction Plato drew near when he put a world
of different Ideas in place of Parmenides' ``One.''  Thereby
not all difficulties were overcome, but better conditions
had been provided for contending with them.  And
despite all the obscurities of the dialogue \textit{Parmenides}, Plato
has here become aware that it is not enough to
distinguish between Ideas and particulars, but that the problem of unity
and multiplicity also arises within the doctrine of Ideas.  The
question how a particular, from falling under one
Idea, can come to fall under another Idea, leads Plato
in his own way to take up Zeno's problem.  That he
acknowledges it as a real problem, whereas Parmenides
and Zeno regard it as the fruit of an illusion,
comes from the world of Ideas being for him a multiplicity.
Each single Idea is for him a reality; therefore the
transition from one Idea to another also had to be a reality.
Zeno laid the presupposition of a real multiplicity at the foundation
only with a polemical intent.  Plato is now, as
we have seen, in agreement with Zeno in this, that from an absolute
multiplicity (a chaotic series of differences) no knowledge
is possible (Parm. p. 164 B--166 E).  But from the very
beginning of his philosophizing Plato had, in agreement with Heraclitus,
the thought that what is given in experience, the particulars, was
caught up in constant becoming and change.  Since to each
group of the particulars caught up in change there now corresponds an Idea,
the problem of the transition of particulars from the domain of one Idea
to the domain of another Idea must necessarily come to
present itself with quite a different force from when one
speaks merely quite hypothetically of motion and change.  And
here the problem of the sudden then springs forth with force. ---

b. Plato has here come upon the concept of the sudden through
the opposition between the sharply distinguished Ideas and the
% --- p. 44 ---
\setcounter{footnote}{0}%
\opage{44}particulars' continuous stream.  But this concept has also emerged at a quite different place
in Plato's philosophy.  Whereas
Parmenides and Zeno were not psychologists,
psychological interest is a very prominent trait in Plato.  It was
not merely dialectic and metaphysics that led him beyond
the Socratic standpoint; it was quite as much
his deep psychological understanding.  His doctrine of Eros (in the
\textit{Symposium} and the \textit{Phaedrus}) has rightly been called the
most important contribution of ancient thinking to psychology.  Whereas
Socrates once for all presupposed that in human beings there
must be both urge and capacity to form thoughts
that could be held fast through all changes, and only strove
to get them to let the urge for clarity
and truth gain influence on their lives, Plato seeks to
show how this urge, the intellectual longing, and
the corresponding capacity can develop step by step.
Eros, which for him in the widest sense designates the whole life of the drives,
appears first as an urge for sensible satisfaction,
but can nevertheless point beyond this.  And the very urge
that stirs after sensible satisfaction can prove to
serve ends that lie beyond the individuals each
taken separately.  Through the urge to procreate the single
individual works, without knowing it, for the continued existence of the race, and in the parental instinct
it appears quite clearly that something stirs here
that leads beyond individual self-preservation.  The
capacity and urge to create and secure life that express themselves in
these drives are awakened by beautiful, attractive objects, which
come to stand as the aim of the individuals' urge.
But from bodily beauty one can pass over to spiritual
beauty as the releasing force, there being a relation of likeness or
kinship (ἀδελφόν, συγγενές, Symp. 210 B C) between
everything that is beautiful.  And by virtue of the same relation
% --- p. 45 ---
\setcounter{footnote}{0}%
\opage{45}one can pass over from practical striving to intellectual
striving, striving after truth for its own sake.
Through a series of steps the individual can thus raise
himself to knowledge, or rather vision, of the highest Ideas. ---

The line of thought underlying this Platonic
conception has its significance to this day.  When
for Plato the urge of love, the urge for honor, the poetic urge, and
the urge for knowledge can form links in a series of development, it
is because they can form a series of what one may
call psychological equivalents.  It is a series of
displacements of motive that he describes in his doctrine of the great
Eros.  There is ideality in every drive, insofar as it points
beyond the given state, beyond what has hitherto been
attained, --- makes a demand.  And in this Plato sees the possibility
of a metamorphosis, in that a thought of an end can, by virtue of
likeness, replace another.  Eros can become the great
teacher, in that the drive, once it has been awakened, can
gain new, perhaps higher ends\footnote[1]{Cf. my treatise on the \textit{Symposium} in the third volume of Mindre Arbejder.}.  If this is correct,
there lies hidden a deep psychological insight in Plato's doctrine
of Eros, and it is something other and more than ``ein geistreiches
Spiel,'' as Wilamowitz-Moellendorff has recently called it\footnote[2]{\textit{Platon,} II, p. 172--174.}.

But whatever one may think of this, it is evident
that in the domain of the life of the drives Plato assumes the possibility
of a psychological process of development, within which
by virtue of definite laws a transition takes place between highly
different stages.  The moments of life here do not stand, as
they would have to on Zeno's presuppositions, with absolute independence
over against one another, but form (or can at least
form) a series in which the later links ever
% --- p. 46 ---
\setcounter{footnote}{0}%
\opage{46}presuppose the preceding ones, so that the links cannot be interchanged;
the order of succession is determined; --- the order of the summands is
not a matter of indifference.  At every point the individual is at rest,
insofar as his urge, as it has developed
through the preceding stages, finds fulfillment;
but he is at the same time in motion, insofar as the step attained
(ἐπαναβαθμός) forms the presupposition for, or contains an
incentive to, going further.  Only when one advances
in a definite sequence (ἐφεξῆς καὶ ὀρθῶς) is
the highest goal a human being can reach attained.  Only he reaches
this goal who by going through the series of stages has
become strengthened and has grown (ῥωσθεὶς καὶ αὖξηθείς) (Symp.
210 D E).  But then at last the highest beauty may suddenly
(ἐξαίφνης) reveal itself to him (210 E).

Whereas the sudden appeared suddenly and surprisingly
from the standpoint of Eleaticism and the doctrine of Ideas, where one went from
being to being, from Idea to Idea, the sudden
in a way does not come suddenly in the \textit{Symposium}'s doctrine of displacement,
which impresses the definite sequence of the
preceding stages, concluded by the highest vision.
A psychologically motivated preparation precedes
the last link in the series, which nevertheless also belongs to
the series.  It is a kind of contrast effect that makes itself felt,
and such an effect has the remarkable feature that it both separates,
namely through the opposition of quality between the links,
and unites, namely in that the quality of the following link is
determined precisely by the opposition to the preceding one.  The transition
from the next-to-last to the last link is determined by the law
that holds for the transition between all links in the series.
However valuable the last link may be in relation to
the preceding ones, --- even if it denotes a state in which
life first truly acquires value for a human being, --- it nevertheless denotes
% --- p. 47 ---
\setcounter{footnote}{0}%
\opage{47}a definite point of time, which has its definite place
within the course of life.  (Symp. p. 211 D: ἐνταῦθα τοῦ βίου,
ἐίπερ που ἄλλοθι, βιωτὸν ἀνθρώπῳ).  What occurs suddenly
comes like a ripe fruit, conditioned by the growth
and flowering of the past.

In his interesting study ``La théorie platonicienne de
l'amour'' (Paris 1908) Léon Robin strongly emphasizes the
complex, synthetic character of the Platonic Eros.  Contrary
oppositions make themselves felt in this feeling or drive,
and Robin rightly thinks that at this point in Plato's thinking
an influence of Heraclitus shows itself; this philosopher indeed laid such great weight on the
contrasts that existence presents.  And in the fact that through
the state, rich in oppositions, that Eros denotes, an intimate
relation to the Ideas can be reached, Robin finds an approximation
toward mitigating the sharp difference between Ideas
and particulars.

Since in this treatise I have come to the doctrine
of Eros on the occasion of the concept of the sudden, it is
quite naturally the successive contrasts in the psychological process
springing from Eros that interest me.  It
is expressly said, after all, that Eros as the son of Poros and Penia
is now rich, now poor, now full of life, now dying, and
it is through such opposite phases that a human being
grows to attain the great moment when the Highest suddenly
reveals itself.  It is, so to speak, through a struggle for
existence that this high point is reached.

In the fact that through psychological development one can advance
from participation in one Idea to participation in another
Idea there already lies a hint that the Ideas
must stand to one another in an inner relation.  That they at the same time
can be united in the same particular already pointed in this
% --- p. 48 ---
\setcounter{footnote}{0}%
\opage{48}direction.  But it appears still more clearly in that they
(or the relations to them) can denote different stages in
a continuous development.  Even if they meet here only to
part at once, Plato himself maintains in his dialectic
that it is both the difference of thoughts and their likeness
that brings them into relation to one another.  Plato's doctrine of Ideas
and his doctrine of Eros point each in its own direction, in that Eros
is unintelligible from the doctrine of Ideas, which contains no
possibility of ascribing to striving and development any
real significance.  But as I have sought to show, the
doctrine of Eros nevertheless contains points of view that can lead
to a remedy for the difficulties of the doctrine of Ideas.  It is the opposition of
knowledge from experience to speculation that makes itself felt here.  And after Plato in the \textit{Parmenides} had so
sharply brought out these difficulties, it was to be expected
that he would also by another way than the psychological
way he had followed poetically in the \textit{Symposium} and the \textit{Phaedrus}
take up the problems that were bound to present themselves here. ---
Before I go on to examine the new attempts at thought
that Plato thus comes upon, I will first bring forward
some hints that lead in a similar direction to the doctrine
of Eros, in that they depict a spiritual development toward the conception of Ideas
as an aim.

Not only in the \textit{Symposium}, but also in the \textit{Phaedo} and in the \textit{Phaedrus}, the sudden is mentioned
in the appearance of the highest spiritual state.  From the lower world in which we
human beings live, it is nevertheless possible, according to the \textit{Phaedo} (p. 109 E),
now and then to surface and catch a glimpse of
the world of Ideas, just as the fishes now and then can surface
from the sea and see something of our world.  In the \textit{Phaedrus}
(p. 248 A) the charioteer of the human chariot of the soul,
in the course of the journey in pre-existence, just manages to get his head up
% --- p. 49 ---
\setcounter{footnote}{0}%
\opage{49}over his accustomed horizon, so that he can catch a glimpse of true
reality.  In these passages the contrast between the highest stage of human
existence and the preceding stages, which are nonetheless its necessary
presuppositions, is expressed in pictorial form.  As the doctrine of Eros shows,
it is not a purely theoretical development that reaches its goal
at that highest stage.  What is presupposed is the development of the whole personality,
an appropriation along the way of life.  Already in the \textit{Symposium}
this is stated, where it is said that there takes place not merely a
vision, but also a living together with the Highest (θεᾶσθαι
καὶ συνεῖναι p. 212 A).  And it is stated still more definitely in
the seventh book of the \textit{Republic} and in the Seventh Letter.  There it is taught
that the Highest is not knowledge, but a relation to something
that is raised above all knowledge, and to which we therefore also
cannot ascribe existence in the ordinary sense of the word.
Although, according to Plato, it is possible only for one who has
passed through a rigorous scientific schooling to reach this
point, it is nevertheless a personal experience that in the end
is decisive, so that only an indirect and pictorial communication
to others is possible, and then only after one has been turned
not merely with thought, but with the whole soul, in
the direction of the Highest (\textit{Republic} VII p. 518 C D).  In the Seventh
Letter (p. 341 C; 344 B) both the
sudden and the personal experience are sharply emphasized.  As a protest
against the young Dionysius's priding himself on being initiated
into Plato's deepest and most inward thoughts, Plato declares
here: ``If anyone now or later claims to possess
a knowledge of that which my striving really concerns, whether
they may have heard it from me or from others,
or may have discovered it themselves, then I declare that it is
impossible that they can know anything about it.  I have written nothing
about it and shall never write anything about it, for
% --- p. 50 ---
\setcounter{footnote}{0}%
\opage{50}it cannot be communicated like other kinds of knowledge.  But after
work upon it and after one has taken it up into one's life,
a fire suddenly (ἐξαίφνης) blazes up, as if a
spark had fallen into the soul.  Once this fire is kindled,
it feeds itself''\footnote[1]{What is emphasized in this statement, the sudden and the unutterable or only indirectly communicable, is also impressed by \textsc{Plotinus}.  See \textit{Ennead.} V, 3, 17 (τότε δὲ χρὴ ἑωρακέναι πιστεύειν, ὅταν ἥ ψυχὴ ἐξαίφνης φῶς λάβῃ) and VI, 9, 4 (--- κατὰ παρουσίαν ἐπιστήμης) κρείττονα . . . οὖδὲ ῥητὸν οὖδὲ γραπτόν).}.

Here too the sudden appearance of the new state is
an essential feature, together with the fact that spiritual growth and personal
immersion are conditions for it.  --- Psychologically speaking,
sudden experiencing and long preparation do not exclude each other.

In the doctrine of Eros emphasis is laid on all the powers of the soul
working together in a concentrated manner.  But even
where Plato is chiefly occupied with the purely intellectual
development, there is for him no contradiction in the
highest stage, in spite of its contrast to the preceding
stages, being nevertheless prepared by them.

In the \textit{Phaedo} (p. 74—79) the development of strictly scientific
knowledge is depicted, starting from the connection between
representations that is brought about by the association of representations
through likeness and by involuntary comparison.  Sensible
things, such as pieces of wood and stones, resemble one another in
certain respects, not in others, and perfect likeness (identity)
does not occur in sensible experience.  And yet it is by
occupying ourselves with those imperfect likenesses that we work our way
forward to the Idea of likeness!  And as it goes with likeness, so it goes
with other thoughts, with the thoughts of justice,
beauty and goodness.  In what experience shows us there is at
most a tendency, a striving toward what the Idea contains,
but experience can nevertheless serve as a preparation
% --- p. 51 ---
\setcounter{footnote}{0}%
\opage{51}for the perfected acquisition of the content of the Idea, which only the
highest science bestows.  That this is taught in the \textit{Phaedo} is all the more
interesting because this dialogue otherwise stands in a certain
contrast to the \textit{Symposium} by stressing how great
hindrances to true knowledge sensation and sensible
impulse can be.  Experience for Plato never contains the
full explanation of true science.  He does not know
induction in the modern sense of the word, and does not use
the corresponding Greek word either (ἐπαγωγή), which first occurs
in Aristotle.  What takes place in the course of the intellectual
development is for Plato not an introduction to the
connection of the particulars, but ``an elevation (ἐπαναγωγή)
of the best in our soul to the vision of the best in the true
being'' (\textit{Republic} VII p. 532 C).  It is, as has been aptly
said, a kind of ascension that Plato here describes\footnote[1]{\textsc{Teichmüller}: \textit{Studien zur Geschichte der Begriffe,} p. 407; 425.  Teichmüller conjectures that in the Platonic Epanagoge (elevation) lies the germ of the Aristotelian Epagoge (induction), which is an introduction to the world of the general.  With Aristotle the ascension becomes an introduction to the general laws of the particulars.}.

In a more methodical form Plato has, in the sixth book of the \textit{Republic},
depicted the development of knowledge through a series of stages.
From sensation one ascends to opinion and conjecture,
from there to strict mathematical knowledge and finally
to the testing of the presuppositions of this knowledge, which
according to Plato leads beyond everything hypothetical to knowledge of
what is in itself.  There takes place here a progression from
stage to stage, in which no stage can be dispensed with if one wishes to
reach the highest.  ---

Plato has thus under various forms maintained that,
however high above all sensations, impulses, hypotheses, and
inferences the highest knowledge lies, both as regards
clarity and as regards value, it is nevertheless reached only
% --- p. 52 ---
\setcounter{footnote}{0}%
\opage{52}through faithful striving and work on the lower stages.
But when the question is raised how such a
course of development, such an ascent, is possible at all,
then it indeed stands for Plato as a wonder, which
can be explained only by our souls having already in pre-existence
(even if only in glimpses) stood face
to face with that which must now be acquired amid privation and
longing and through earnest work.  This mythological
explanation sprang from wonder that the
result reached along the long ways seems far to surpass
what was prepared on the way.  Plato's way out here can
recall the solution of the problem of knowledge by modern philosophy of evolution,
only that there one points back from the
experiences of the individuals to the experiences of the race, whereas
Plato points back to the very pre-existence of the individuals themselves\footnote[1]{Cf. my treatise: \textit{Nogle Bemærkninger om Platon’s Psykologi.} (Tidsskrift for Filologi og Pædagogik, 1879), p. 209---213.}.
---

Before I conclude this section, in which the doctrine of Eros has
played an essential role, it must be remarked that this
doctrine does not exhaust the whole of Plato's psychology.  One must,
as I already set forth in my doctoral dissertation ``Den antike
Opfattelse af Menneskets Villie'' (1870), distinguish three
psychological types in Plato.  The first is the one that I have in this
section brought forward especially, and which prevails in the \textit{Symposium}
and in the \textit{Phaedrus}.  It is remarkable in taking its starting point
in elementary natural want and yet leading up toward an
absolute culmination.  The movement proceeds along one line.
The second psychological type that Plato depicts presents
a development from different starting points, in that there is assumed
a difference, respectively a conflict, between three different
tendencies in the soul, --- toward enjoyment, toward self-assertion
% --- p. 53 ---
\setcounter{footnote}{0}%
\opage{53}in spirit and anger, and toward knowledge.  What in Eros welled up
with united force, acts here sporadically, and the task
becomes that of reaching a harmony between the conflicting tendencies.
In the fourth book of the \textit{Republic} this conception is taken as
a basis; it is nevertheless also traced in the \textit{Phaedrus}.  In the third
type all energy is presupposed to be concentrated in the life of thought, which
rules over a power that holds the craving for enjoyment and for self-assertion
down.  In the \textit{Phaedo} this type prevails; but it
also appears in the ninth book of the \textit{Republic}.

The difference between the three types can, as is seen especially
from the investigations of Dr. Hans Ræder, play a role in
the discussion of the chronology of the dialogues.  But they pass through
various metamorphoses into one another.  It is a testimony
to the richness of Plato's psychological conception
that he has brought them all forward; but their opposition
poses great problems.  And what in this connection
especially interests us is what I already in the work mentioned
emphasized\footnote[1]{\textit{Den antike Opfattelse af Menneskets Villie,} p. 64; 74.}, that neither the yearning called
Eros nor the development toward harmony which the second type
conditions can be derived from the doctrine of Ideas itself.  The doctrine of Ideas presupposes
a life of the soul that is once and for all at the height.
How it has got up there it cannot explain, ---
and, as we have seen, and as Parmenides' third objection
establishes, neither how the knowledge that can be
won up there can find real application to the world
of particulars.

\chaphead{4. Later Dialogues in Relation to the Problems\\ of the \textit{Parmenides}.}


The self-criticism which Plato, according to the conception maintained
in the foregoing, exercised in the \textit{Parmenides}, has left
% --- p. 54 ---
\setcounter{footnote}{0}%
\opage{54}effects in his later thinking, although he did not
refute the objections to the doctrine of Ideas set forth in the \textit{Parmenides}.
His later doctrine cannot be regarded as a
direct continuation of the doctrine of the \textit{Symposium}, the \textit{Phaedo}
and the \textit{Republic}, and yet he has by no means given up the doctrine of Ideas,
although it recedes somewhat into the background.  In order to
be able to ``do something about philosophy'' (Parm. 135 C), that is, in order
that his thinking should not come to a standstill, he had to adopt a different
formulation of problems from the one he had hitherto employed.  Such
a new formulation of problems is traced in the whole group of
dialogues that lies between the \textit{Parmenides} and the \textit{Laws},
the latter dialogue not treating questions that concern
the doctrine of Ideas.  This group consists of the \textit{Theaetetus}, the \textit{Sophist},
the \textit{Statesman}, the \textit{Philebus}, and the \textit{Timaeus}.  Recent scholars
(Jackson, Brochard, Ræder, Robin) agree that
there is an inner connection among these dialogues
with respect to problems and line of thought, even if
the chronological relation between them is conceived differently.

I cannot follow Dr. Ræder in placing the \textit{Theaetetus}
before the \textit{Parmenides}, since the way in which it takes up
the problem of knowledge is intelligible only against the background
of the \textit{Parmenides}' thorough critique of the doctrine of Ideas, and since the treatment of problems
in the \textit{Theaetetus}, as I shall attempt to
show, is a step in the direction that is carried further in the \textit{Sophist}
and the \textit{Statesman} and concluded in the \textit{Philebus} and the \textit{Timaeus}.
Therefore there may well be, as linguistic grounds are thought
to argue, some interval of time between the \textit{Theaetetus}
and the \textit{Sophist}.

Already in the \textit{Republic} and in the \textit{Phaedrus} knowledge was described
as an interaction of connection and distinction,
seeing together (synopsis) and analysis, and in this interaction
there came to light what Plato called dialectic.
% --- p. 55 ---
\setcounter{footnote}{0}%
\opage{55}From this it is seen that the abstract developments about unity and multiplicity
in the second section of the \textit{Parmenides} had by no means
been necessary for Plato to be able to continue his
philosophizing, after he had acknowledged the objections
to the doctrine of Ideas in its original form.  The very concept
of thinking as a constant interaction of connection
and distinction, which he had reached in the two dialogues mentioned,
would have been a sufficient starting point, and it
is in fact this starting point, not those abstract
speculations (or sophisms), that he uses in the following
dialogues, which mark a significant change in his
philosophy.

In the \textit{Theaetetus} it is maintained, --- in criticism of the claim that
knowledge should consist in passive and momentary sensation,
--- that all knowledge presupposes a synthesis, without which
all distinction, all analysis would be impossible.  The various
sensations do not sit inside us like the warriors in the Trojan
horse, but are united in a comprehensive unity,
whether one will call this a soul or give it
another name (p. 184 D).  In the very fact that
sense perceptions can be distinguished from one another and come
to stand as equal or different, a
synthesizing function comes to light (p. 185 A).  This consideration
then leads on (p. 185 C) to setting up being
(position) and non-being (negation), identity and difference,
as well as the concept of number as what one might call, with a modern expression,
fundamental categories, i.e. concepts with which we
work in all our thinking, no matter what it is
we are thinking about.  That such a special group
of concepts or Ideas is brought forward in close connection with the very
activity of knowing marks a continuation of the
% --- p. 56 ---
\setcounter{footnote}{0}%
\opage{56}distinction that was made in the \textit{Parmenides}, more than in earlier dialogues,
between different groups of Ideas.

In the \textit{Sophist} Plato goes more closely into the activity of thought itself,
in that he takes the significant step of
investigating judgment as an expression of thinking.  In every
judgment different concepts are connected.  Judgment is an interweaving
of Ideas (διὰ γὰρ τὴν ἀλλήλων τῶν εἰδῶν συμπλοκὴν
ὁ λόγος γέγονεν ἡμῖν, p. 259 E).  This very connecting
thinking, it is said further, is something that is, something that exists,
and with the dropping away of which all philosophy would fall away
(260 A).  Possibly there lies here an answer to the question
that at the critical point in the \textit{Parmenides} was put to
the young Platonic Socrates: ``what will you now do about
philosophy?'' --- and an answer that fully outweighs the
series of pseudo-dialectical considerations which make up the second
part of the \textit{Parmenides}.  It had already in the \textit{Parmenides} itself
been maintained against the Eleatics that there was no
contradiction in the same particular's being able to fall under different
Ideas; already this points to a certain connection
among the Ideas themselves.  And after the \textit{Theaetetus} has
shown synthesis to be the presupposition of the mutual
relation of sensations as equal or different, there is now
undertaken in the \textit{Sophist}, through an analysis of judgment, a consideration
that shows that the different fundamental concepts or categories
can no more than the sensations be absolutely separate
and independent of one another.  Non-being can thus very well
be a predicate of a being, when it is a matter of designating
the latter's difference from another being (ὄντος πρὸς ὄν ἀντίθεσις
257 E).  Likewise one and the same being can stand as like
and different, as at rest or in motion (256 B).  To the categories
mentioned in the \textit{Theaetetus} there are in fact added here rest and
motion.

% --- p. 57 ---
\setcounter{footnote}{0}%
\opage{57}In his doctrine of Eros Plato had already, without being
aware of its epistemological significance, given
an example of such an interweaving as every judgment
means.  For in the constant striving or tendency there is
both identity (as regards the goal) and difference (by reason of
the distance from the goal), both rest (in the constant direction)
and motion (in the restless yearning).  And in the concept of
soul that is set up in the \textit{Phaedrus} (p. 245 C) as that which moves
itself, and which therefore must be in eternal motion\footnote[1]{The concept αὐτοκίνητον seems to derive from Alcmaeon.  See \textsc{Diels}: \textit{Vorsokratiker}’, p. 105.}, there is
also presupposed rest and motion at once. ---

The circumstance that Plato in the \textit{Sophist} takes
an analysis of the function of judgment as a basis recalls Kant's
procedure in the ``Kritik der reinen Vernunft'', and modern
Kantians have accordingly also found critical
philosophy outright in the last period of Plato's philosophizing.  Now
Plato and Kant do indeed meet at this point, and it
is a meeting of world-historical significance.  But they meet
in that Plato ends where Kant begins.  And a pure epistemologist
Plato never became.  The problem of knowledge
did not become for him properly the fundamental problem, or rather,
it was for him solved beforehand in his metaphysics.
He holds fast, in spite of everything, to his doctrine of Ideas.  A philosopher is and
remains for him one who through his thinking relates
himself to the Idea of what is (ὁ δὲ γε φιλόσοφος, τῇ τοῦ
ὄντος ἀεὶ διὰ λογισμῶν προσκείμενος ἰδέᾳ. Soph. 254 A).
That with his analyses in the \textit{Theaetetus} and the \textit{Sophist}
he must consistently come out positively beyond his own doctrine of Ideas,
just as he had come out negatively beyond it in the
\textit{Parmenides}, he does not notice.  There did not present itself to
him the problem which, on the basis of similar analyses
% --- p. 58 ---
\setcounter{footnote}{0}%
\opage{58}c. 2000 years later presented itself to Kant: if our categories
are ultimately forms of our own activity of thought,
whence comes then their validity for objects that
are not produced by our activity of thought?\footnote[1]{See \textit{Kontinuiteten i Kants filosofiske Udviklingsgang} p. 52.}  Even where
Plato finds ever so much self-activity and ever so much
unrest in our thought, the heaven of Ideas is not touched
thereby for him.  The question how he then knows that
there is such a heaven he does not raise.

Bonitz, in his excellent analysis of the \textit{Sophist}, has already
shown\footnote[2]{\textit{Platonische Studien?} p. 195 f.} that Plato was not aware that he
was on the very point of dissolving all metaphysics into logic.  One can find
in him a first possibility of the significant turn
that occurred with Kant's definitive foundation of
critical philosophy.  But this possibility Plato did not see
and did not use, --- was neither able nor willing to use
it.  He remained in agreement with the Eleatics that every
true thought must be a thought of something that has its
absolute subsistence (Parm. p. 132 B; 135 B). ---

The dialogue \textit{Statesman} has, as the name indicates, a
practical task.  But in order to find the ``royal science''
(p. 292 E) or art (p. 300 E) that must be demanded
of the true statesman, a circumstantial division
of arts and sciences is undertaken, a division that at the same time --- like
a corresponding classification in the \textit{Sophist}, --- supplies
examples of how the place of concepts in the world of thought
is determined by their mutual likeness and difference.  And
on a given occasion the importance is emphasized of a standard of measure
for finding the right thing in the individual case.  The
practical art or science presupposes not merely the ability
to compare magnitudes in general, but especially the ability
% --- p. 59 ---
\setcounter{footnote}{0}%
\opage{59}to find the determinate quantitative proportion between different
capacities and impulses, which must be present if the tasks
that ethics and political theory set are to be satisfied.
It holds for all practical endeavors that they can only
lead to what is beautiful and good when the right measure
is observed (τὸ μέτριον σώζουσαι πάντ’ ἀγαθὰ καὶ καλὰ ἀπεργάζονται.
Statesman p. 284 B). --- The ethical significance of the art of measurement
had already been emphasized in the \textit{Protagoras} (p. 356—357)
and in the \textit{Gorgias} (p. 507—508).

This last consideration is carried further in the \textit{Philebus}, which,
in connection with the \textit{Timaeus}, gives the epistemological
conclusion of the movement of thought that began in the
\textit{Parmenides}, and thereby stands as the decisive expression
of Plato's later philosophy.  Some authors see in this
dialogue the \textit{Philosophos} that Plato in the \textit{Sophist} (p. 217 A)
had announced as the third member of a trilogy, whose first two
members were to be the \textit{Sophist} and the \textit{Statesman}.

In the \textit{Sophist} an inner connection between the Ideas
was established by the fact that the basic form of all activity of thought, judgment,
is a connection of concepts, a ``mixture.''  In
the \textit{Statesman} the necessity was emphasized of a determinate
quantitative proportion, a determinate proportion of measure, when it is a matter of the decision
of practical value.  In the \textit{Philebus} it is now shown how
in all understanding and in all evaluation it is a matter of
finding a determinate relation between a limit (rule) and an
unlimited, which is to be limited.  By the unlimited
(τὸ ἄπειρον) is thought especially of the incoherent diversities
(cf. Statesman p. 273 D: ὁ τῆς ἀνομοιότητος τόπος
ἄπειρος ὤν).  Plato's struggle for his Ideas is constantly a
struggle against chaotic series of differences.  Neither in the world of Ideas
nor in that of particulars can such series of differences be tolerated.
Everything that is consists of unity and multiplicity and
% --- p. 60 ---
\setcounter{footnote}{0}%
\opage{60}has limit and unlimitedness in its nature.  (Philebus p. 16 C).
The limit comes to light in that identity, number, and
measure rule (25 A).  The dialogue's proper subject is psychological
= ethical, therefore the standard of measure for evaluation plays a
greater role than the standard of measure for understanding, although there
are common presuppositions for both.  Whereas mathematics in
the \textit{Republic} is conceived as a stage of knowledge that rests
on certain presuppositions, whose ground only dialectic
reveals, it furnishes in the \textit{Philebus} means and form for the determination
of the essence of things, quite especially for the determination
of values (``the Good'').  The concept of the right measure
(τὸ μέτριον) now becomes a central concept.  In that ``the
Good'' presupposes a proportion determined by measure and harmonious
(μετριότης καὶ συμμετρία), it appears as beautiful (p. 64 E).

Plato here involuntarily comes to hint at the connection
between the concept of value and the concept of totality.  For
if it is asked what then results from the right measure
not being maintained, the answer is that it leads to dissolution into an
unhappy chaos (p. 64 DE cf. Statesman p. 273 D).  In our
modern language we can express this by saying that every whole
demands certain inner and outer conditions for its subsistence,
and that the evaluation we make of the states and actions of personal and
social wholes is ultimately grounded
in the intervention of these states and actions in
the conditions for the subsistence and development of the whole concerned\footnote[1]{Cf. \textit{Totalitet som Kategori,} p. 66---73.}.
The possibility of a rational ethics rests on the possibility
of exact determination of such conditions of life.
In the \textit{Republic} Plato has sought to show that the individual personality
and the conditions of life of the state stand in a harmonious relation
to each other, when both wholes reach their perfect
development.  In his later years he seems to have laid the whole
% --- p. 61 ---
\setcounter{footnote}{0}%
\opage{61}weight on the quantitative relations that must be fulfilled
in order for these conditions of life to be present.  It is reported
that those who came to hear Plato's lecture on the concept of
``the Good'' were to a high degree disappointed at hearing only
talk of arithmetic and astronomy, of ``the limit'' and of ``the
One.''  On the other hand he has in his statements cited above
in the seventh book of the \textit{Republic} and in the Seventh Letter
raised ``the Idea of the Good'' so high above all knowledge that
he thereby deprives himself of the rational foundation for ethics and
at the same time tears asunder the bond between the Good and the
Beautiful that stands firm for him as long as he moves within
human experience.  The mystical and the rational
element in his philosophy here enter into sharp opposition to
each other. ---

No more than the epistemological line of thought in
the \textit{Sophist} does the doctrine of limitation and of the necessity
of determinateness of measure for every whole (``mixture'')
lead Plato to give up the doctrine of Ideas.  The highest knowledge is still
that of ``what truly is, that which always conducts
itself in the same way'' (p. 58 A).  The doctrine of Ideas remains in the background,
after it had proved necessary to acknowledge
the connection which the \textit{Sophist} has demonstrated through analysis
of judgment, and the character of mixture (quality of wholeness)
which the \textit{Philebus} has demonstrated in everything that subsists and has value. ---

In the \textit{Timaeus} it is a main thought that the connection of
identity and difference is expressed in the structure of the universe.  In
the celestial equator identity appears through the uniform movement
of the celestial sphere; in the ecliptic the diversities appear
through the movements of the sun and the planets.
Corresponding to this, identity comes to light in man's
thinking, diversity in sensation.  The same
opposition makes itself felt in the universe-encompassing
% --- p. 62 ---
\setcounter{footnote}{0}%
\opage{62}world soul, which conditions the mathematical lawfulness of nature.
This connection of identity and difference corresponds to
the connection of concepts in judgment in the \textit{Sophist} and to the
connection of limit and unlimited in the \textit{Philebus}.  It was the
new doctrine presented in the dialogues mentioned that made it
possible for Plato to obtain a framework within which he
could order the physical, astronomical, and biological material
that he collected in his later years.  It has rightly been remarked
that it would be a misunderstanding to believe that in
the Academy only philosophy, mathematics, and astronomy were cultivated.
Biological studies also occupied minds, and perhaps
the foundation was already laid here for Aristotle's interest in
biology\footnote[1]{See \textsc{Burnét}: \textit{Greek Philosophy. I. Thales to Platon,} p. 223.}. % print: Burnét (misprint)
While in his classical doctrine of Ideas (in the \textit{Symposium},
the \textit{Phaedrus}, the \textit{Republic}) he was content to ascend
from particulars to Ideas and to go from Idea to Idea within the
world of Ideas, he had now put forward thoughts that made it
possible for him to attribute to the particulars another significance than
that of merely being a springboard on the journey up to the world of Ideas,
and it became a task for him to find a connection
among the particulars themselves.  Because of the imperfect
means available, an attempt in this direction
was bound to appear childlike, and Plato makes ample
use of mythological representations.

Both from the philosophical and from the philological side it has been shown
that in the \textit{Timaeus} an influence of Democritus's
thoughts makes itself felt, especially of his atomic doctrine.  The
strict lawfulness which Plato maintains (although he
subordinates it to a teleological point of view) holds
not only for the world of sensible particulars, but also for the
elements, in themselves non-sensible, of which it ultimately consists.  Already
Democritus's atoms were non-sensible, and already in antiquity
% --- p. 63 ---
\setcounter{footnote}{0}%
\opage{63}attention was paid to the likeness between Plato and
Democritus, that both, each in his own way, explained the
sensible qualities by non-sensible causes.  Plato goes
further back than Democritus, in that the atoms according to him
must be assumed to be composed of elements that can only
be characterized purely mathematically, and at the basis of these
elements there lie again (as a kind of atoms of the third degree)
finer elements, which are ``known only to God and to him whom he
loves'' (p. 53 D).  In an interesting manner the relativity of the concept of the atom
is here hinted at in mythological form\footnote[1]{\textsc{Léon Robin} (\textit{La Place de la Physique dans la Philosophie de Platon,} Paris 1919), who has strongly drawn attention to the passage cited in the \textit{Timaeus}, seeks to show that by the atoms of the third order one must think of the ideal numbers and ideal figures with which Plato (according to Aristotle) became so much preoccupied in his last, Pythagoreanizing period. --- Through Plato's natural philosophy there runs, alongside much that is fantastic, a striving to trace all qualities back to purely quantitative relations --- although this too happens in a fantastic manner.}.

In thus introducing a third world into his
philosophy alongside the world of Ideas and the world of sensible
particulars, he undertakes in accordance with this a
remarkable change in his theory of knowledge.  The atoms can
neither be thought (in the manner in which the Ideas can be
thought) nor be sensed (like the particulars with their qualities).
Only the way of analogy then stands open.  It is this that Plato
points to when he says that that which lies at the basis of
the sensible qualities must be grasped ``without the help of the senses
by a kind of spurious thinking'' (p. 52 B).  (Even at p. 27 D
only thinking and sensation are mentioned, without any hint of a
third kind of knowledge)\footnote[2]{See my treatise \textit{Platon og Demokritos} (Mindre Arbejder. Tredie Række), p. 99.}.

It is not certain, however, that the triad of Ideas, sensible
particulars, and atoms of different degrees, which has
% --- p. 64 ---
\setcounter{footnote}{0}%
\opage{64}called forth a corresponding triad within the theory of knowledge,
is due merely to the influence of Democritus's doctrine.  Léon
Robin, who relies especially on the \textit{Timaeus} p. 52 C, has (loc.
cit. p. 26) rightly emphasized that, if the particulars could not
be derived from the Ideas, a factor had to be presupposed
which could make it possible that the particulars, in spite of their
likeness to the Ideas, could nevertheless have properties in which they
deviated from them.  There must be something by which the image is
different from the model; --- there must accordingly be given something
different from the Ideas (ἕτερον τι p. 52 C).  If there
were given nothing but the Ideas, their images would be inexplicable.
It is by a postulate, which is grounded in the difference
between Ideas and particulars, that Plato from his earlier
doctrine arrives at the principles of his natural philosophy in
the \textit{Timaeus}.  He here in fact draws the consequence of the
third objection in the \textit{Parmenides} against the doctrine of Ideas, --- a
testimony that he felt himself hit by it. ---

If there are in the line of thought that Plato enters upon in the
\textit{Timaeus} features that make it intelligible that the natural
scientists of the Renaissance, especially Galileo, could find a kindred spirit in
him, this is understood still more clearly through the formulation of problems
that Plato, according to Simplicius, is said to have assigned to astronomy.
He finds its task, namely, in investigating
what regular motions must be assumed in the universe,
in order that the phenomena present in experience may
come into their own.  It is a way of asking that
presumably stems from the Pythagoreans, but Plato gave it
a stricter form, and Plato's friend Eudoxus sought to solve
the task that it set.  As Duhem, who has especially
investigated this point, says: ``The method of the physical sciences
has been defined by Plato and by the Pythagoreans
of his time with a clarity and an exactness that
% --- p. 65 ---
\setcounter{footnote}{0}%
\opage{65}have not been surpassed''\footnote[1]{\textit{Le Système du Monde. I. De Platon à Copernik.} 1913, p. 129.}.  If this is correct, Plato has
through his epistemological demands upon science
exercised an indirect influence on the development of natural science,
just as according to Zeuthen he has exercised such an
influence on the development of mathematics. ---

It is a great question whether Plato would have come upon
the significant epistemological Ideas
that he has developed in his later dialogues, Ideas that give
these dialogues a great interest for the history of thought
in spite of their artistic imperfection, if he had not
with such great force and love of truth thought through
the objections that could be raised, and to some extent have been
raised, against his doctrine of Ideas.  They did not lead him to give up
the latter, but they have acted as a spur to renewed intellectual work
and have thereby contributed to giving his thinking
a greater universality than it would have had if
it had dogmatically closed itself around the original doctrine of Ideas
as the complete expression of the truth.

% PDF 68, the print's Indhold (p. [66])
\clearpage
\begin{center}{\large C O N T E N T S}\\[0.3em]\rule{3em}{0.4pt}\end{center}
\hfill{\small Page}\par
\tocline{}{Introduction}{1}
\tocline{1.}{Critique of the Doctrine of Ideas}{6}
\tocline{2.}{The Second Part of the Dialogue}{26}
\tocline{3.}{The Concept of the Sudden in Plato's Philosophy}{38}
\tocline{4.}{Later Dialogues in Relation to the Problems of the \textit{Parmenides}}{53}
\begin{center}\rule{3em}{0.4pt}\\[-0.8em]\rule{3em}{0.4pt}\end{center}
\vfill
\begin{center}{\small Presented at the meeting of 16 January 1920.\\ Printed 15 June 1920.}\end{center}
\end{document}
